%% 第24回編集委員会（物理第3巻）：組 波W7
%% p3-13.tex  波-13　光の波：干渉・回折・偏光・散乱（★3〜5・3段）
%% 中身は計画書 tools/phys3-plan.md の B「波-13」。公式の持ち主は A（A38〜A43・A59・A60・A63）。
%% 数値はすべて Python で計算した：checks/物理第3巻/波-13/calc_p3_13.py（出力 calc_p3_13_out.txt）、
%%   coords_p3_13.py（グラフの点と③のプログラムの出力。出力 coords_p3_13_out.txt）。
%% 記号（この章。巻全体は計画書 C-1）：
%%   光の速さ（真空）c、物質の中の光の速さ v、屈折率 n、波長 λ（膜の中の波長 λ/n）、振動数 f、光の強さ I（W/m²。この章には慣性モーメントが出ない）、
%%   スリットの間かく d（回折格子の溝の間かくも d）、スリットのはば a、円い穴の直径 D、スクリーンまでの距離 L、しまの間かく Δx、
%%   干渉の次数 m（整数）。電子の質量は m_e と書き、次数の m とぶつけない。膜の厚さ t、屈折角 r、入射角 i、
%%   時刻 t は②の sinωt と⑧の cosωt・d²x/dt² だけに出る。膜の厚さ t（④）とは別の量なので、②と⑧の本文で断った（修3-6、P6-5）。
%%   くさびのすきまの高さ（髪の毛の太さ）H、くさびの長さ L、レンズの曲率半径 R、輪の半径 ρ_m（屈折角 r とぶつけないため ρ。修3-6、P6-5）、
%%   回折格子のすきまの本数 N（③）、干渉計で動いたしまの本数 M（⑦。N とぶつけない）、⑤のフーリエ変換の変数 q＝2π sinθ/λ（9-4 の ω にあたる。光の角振動数 ω とぶつけない）、
%%   偏光板の向きの角 θ、ブルースター角 θ_B、光のゆれ（電場）の振幅 E_0、電子のゆれの振幅 A、電子をゆする力の振幅 F_0、
%%   電子の固有角振動数 ω_n（6-13・波-2 と同じ）、光の角振動数 ω、散らされる割合の目安 τ（光学的な厚さ）、空気の量の倍率 X。
%%   第4部なので c は光速（W-4）。減衰係数 c は使わない（電子のモデルは減衰を無視する）。波数 k は使わず 2π/λ と書く（ばね定数も出さない）。
%%   熱-13 は比熱とぶつけないため光速を c_0 と書いた。この章では c（思い出そうで一言）。
%% 公式の持ち主：ヤングの実験 Δx＝Lλ/d（A38）、回折格子 d sinθ＝mλ（A39）、薄い膜の干渉 2nt cos r（A40）、くさび・ニュートンリング（A41）、
%%   偏光（A42）、散乱（A43）、円い穴の回折 1.22λ/D（A59）、マリュスの法則とブルースター角（A60）、λ⁻⁴（A63）はこの章。
%%   2つの波源の干渉（A30）は波-8 ③、単スリット sinθ＝λ/a（A29）は波-8 ④、屈折率 n＝c/v（A33）・屈折の法則は波-12 ③（持ち主は波-8 ①）、
%%   反射係数の符号は波-7 ⑤、ガラスの反射率 4 % は波-7 ⑨、強制振動の振幅は波-2 ④⑤、c＝fλ は波-5 ②。まだない章なので文字で指す。
%%   「出す光の強さは加速度の2乗に比例」は電-15 の結果を使う（W-8）。光が電磁波であることも電-15。光の粒は原-2。
%%   熱-13 ①の予告「光の波長と色、光が波であることの詳しい話は、物理シリーズの『波動』（波-12・13）で」は、①（波長と色）・②（波である証拠）・③（波長を測る）で受けた。
%%   力-8 のコラムの約束「ヤングは、光が波であることを示す実験（波-13）でも知られる」は②で受けた。
%% 出典（第24回の追記 2章13：11月5日に確かめた。作業環境から本文を読めたのは physics.nist.gov・arxiv.org（curl）。
%%   WebFetch は www.nasa.gov にも届かず、science.nasa.gov・www.jpl.nasa.gov・emtoolbox.nist.gov・oceanopticsbook.info・JIS の webdesk の本文も読めなかった。届かない資料は WebSearch で、一次資料を引く2つ以上の資料を照らし合わせた）
%%   ・可視光 約380〜780 nm：CIE S 017:2020（e-ILV 17-21-003 visible radiation）＝IEC 60050-845-21-003 は「下の限界 360〜400 nm、
%%     上の限界 760〜830 nm」とし、決まった境はないと注記（cie.co.at/eilvterm/17-21-003、electropedia.org、NIST OSAC の用語集）。
%%     380〜780 nm はこの中の目安。色の境目の波長も、決まった値のない目安（同じ注記）。JIS Z 8120 の本文は webdesk で開けず読めなかった。
%%   ・ヤング："On the theory of light and colours" 1801年11月12日の講演（Phil. Trans. 1802）、"Experiments and calculations relative to
%%     physical optics" 1803年11月24日の講演（Phil. Trans. 94, 1804, pp.1-16）。本文の「1801年から1803年ごろ、講演と論文」はこれでよい。
%%   ・マリュス：1808年、リュクサンブール宮殿の窓の反射光を方解石ごしに見て発見（1808年12月に科学アカデミーへ報告、1809年に発表）。
%%     本文の「1808年」（発見）と「マリュスの法則 1809年ごろ」（発表）はこれでよい。
%%   ・ブルースター：1815年、Phil. Trans. 105, pp.125-159 "On the laws which regulate the polarisation of light by reflexion from
%%     transparent bodies"（1815年2月11日の手紙。Royal Society の目録 PT/9/9）。
%%   ・エアリー：Trans. Cambridge Phil. Soc. 5, p.283（1835年。1834年11月24日に読まれた）。本文には年を書いていない。
%%   ・レイリー：空の色 Phil. Mag. 41（1871年）"On the light from the sky, its polarization and colour"（ブリタニカも1871年）。
%%     分解能の基準 Phil. Mag. 1879年10月 "Investigations in optics, with special reference to the spectroscope"。
%%   ・マイケルソン：1881年、ポツダムの天体物理観測所の干渉計の実験（Am. J. Sci. 22, 1881）。マイケルソンとモーリー 1887年（Am. J. Sci. 34）。
%%   ・メートルの定義 1983年：NIST "Definition of the meter"（physics.nist.gov/cuu/Units/meter.html を読んで確かめた）。
%%   ・CD・DVD・BD のトラックの間かく 1.6・0.74・0.32 μm、読み取りの光 780・650・405 nm：光ディスクの比較表（NA 0.45・0.6・0.85 とともに）と
%%     米国特許の文書（容量 650 MB・4.7 GB・25 GB とともに）、Samsung の解説で一致。DVD の光は 635〜650 nm と書く資料もある。
%%   ・屈折率：MgF2 の膜 1.38（550 nm。光学部品の会社の資料2つ。結晶は 589 nm で no=1.378・ne=1.390）、
%%     レンズのガラス 1.52（SCHOTT N-BK7 のデータシート nd=1.51680、587.6 nm）、水 1.33（20 ℃、589 nm で 1.3330：Tilton・Taylor の測定、
%%     CRC ハンドブックは IAPWS の式の値）。単層の反射防止膜の残りの反射 約1.3 %、反射が赤紫に見えることも資料と一致。
%%   ・ひとみの直径 約2〜8 mm（明るい所 2〜4 mm、暗い所 4〜8 mm。眼科の解説の値。資料の質は高くないが多くの資料で一致）。
%%     視力1.0：ランドルト環（直径 7.5 mm、太さ・切れ目 1.5 mm）の切れ目を 5 m で見分ける＝視角 約1分（コトバンク：日本大百科全書・
%%     ブリタニカ国際大百科事典・大辞林）。本文の明るい所の 3 mm はこの中。
%%   ・すばる望遠鏡の主鏡 8.2 m（国立天文台ハワイ観測所のリーフレット 2024年版、三菱電機技報 1999）。本文には出てこない。
%%     シーイング：マウナケアの大気だけで中央値 約0.4秒（CFHT・Gemini）、すばるの星の像で約0.6秒、ふつうの観測地では1秒前後。
%%     「約0.5〜1秒」は目安としてよい。本文には出てこない（⑤の章末問題の「空気のゆらぎ」だけ）。
%%   ・LIGO：腕 4 km、2015年9月14日 09:50:45 UTC に GW150914 を検出、ひずみの最大 1.0×10⁻²¹（Abbott ほか、Phys. Rev. Lett. 116, 061102 (2016)）。
%%     のび縮みは 4 km×10⁻²¹＝4×10⁻¹⁸ m で、陽子の直径（約1.7×10⁻¹⁵ m）の約400分の1（本文の「数百分の1」でよい）。
%%     KAGRA：腕 3 km、岐阜県の神岡鉱山の地下 200 m 以上（arXiv:1811.08079 ほか）。
%%   ・空気だけの散乱の光学的な厚さ τ(550 nm)＝0.097：Bodhaine ほか (1999) J. Atmos. Oceanic Technol. 16, 1854 の式（海面・1013.25 hPa・
%%     288.15 K）で 0.0971。同じ式で τ(450)/τ(700)＝6.07、λ^-4.08（本文の「λ^-4.1 ほど」でよい）。
%%     地平線の空気の量 約38倍：Kasten・Young (1989) の式で 37.9（Air mass の解説2つで「地平線で約38」）。
%%     いずれも checks/物理第3巻/波-13/calc_p3_13.py の最後で計算し直した。
%%   ・空気の分子の固有振動の波長 約100 nm：空気の屈折率の Edlén (1966) の式の分母が0になる波長は 88 nm と 160 nm（真空紫外）。モデルの目安としてよい。
%%   ・空気の屈折率：Edlén の式で 633 nm・1気圧の n−1 は 0 ℃で 2.92×10⁻⁴、15 ℃で 2.77×10⁻⁴、20 ℃で 2.72×10⁻⁴
%%     （NIST の計算機をもとにした ref_index の値 1.0002716（20 ℃）とも一致）。もとの「15〜20 ℃で 2.9×10⁻⁴」は 0 ℃の値で、ちがっていた。
%%     チャレンジの問題を「1気圧・20 ℃、しま43本、n−1≈2.7×10⁻⁴、n≈1.00027」に直した（46本→43本）。②の「1.0003 を1とする」はそのまま。
%%   ・ミツバチが空の偏光で方角を知ること：フォン・フリッシュ、1949年（ブリタニカ。Naturwissenschaften 1948・Experientia 1949、polarization.com の解説）。
%%     本文には出てこない。
%%   ・火星の夕焼けが青いこと：探査車キュリオシティが 2015年4月15日（956火星日）にゲール・クレーターで写した。チームのレモンの説明は
%%     「細かいちりの大きさのため、青い光が少しだけ大気を通りやすく、ちりで散らされた青い光は太陽の向きの近くにとどまる。ほかの空は黄色〜だいだい色」
%%     （NASA/JPL の発表を引く Spaceflight Now・Christian Science Monitor・ScienceAlert、NASA SVS 11875）。本文の説明と合う。
%%     NASA のページ（science.nasa.gov・www.jpl.nasa.gov）は作業環境から届かず、本文は直接読めなかった。
%%   ・まだわかっていないこと（シャボン膜の寿命）：膜が薄くなる流れ（排液）と蒸発・ちりなどで割れる時間が大きくばらつき、
%%     決まった寿命を予測する理論はまだない、という扱い。論文名を確かめた：L. Champougny ほか "Influence of evaporation on soap film rupture"、
%%     Langmuir (2018) doi:10.1021/acs.langmuir.7b04235（arXiv:1803.00801。arxiv.org で読んだ。蒸発による薄くなり方が排液と同じくらいになると割れる）。
%%     もとのコメントの「J. Fluid Mech.」はまちがいで、直した。割れるまでの時間が大きくばらつくこと：Gilet ほか、APS DFD 2007 "Lifetime expectancy
%%     for a soap bubble"、T. Aurégan・L. Deike arXiv:2411.03908（2024。排液は湿度によらず計算できるが、寿命は湿度で10倍以上変わり、
%%     まわりの空気のかき混ぜで分布が変わる。arxiv.org で読んだ）。本文の書き方（まだわかっていないこと）でよい。
%% 漢字：常用漢字だけ（縞→しま、蝶→チョウ、楕円→だ円、塵→ちり）。checks/物理第3巻/波-13/kanji_check.py で確かめた。
\makeatletter
% 小さな太陽（光線つき）  \rxnsun{(x,y)}{半径}
\newcommand\rxnsun[2]{\begin{scope}[shift={#1}]
  \foreach \a in {0,45,...,315}{\draw[bkOrange!85!black, line width=0.7pt] (\a:#2*1.35) -- (\a:#2*1.8);}
  \filldraw[fill=bkYellow!70, draw=bkOrange!85!black, line width=0.6pt] (0,0) circle (#2);
\end{scope}}
% 立っている人（小さい）  \rxnman{(x,y)}{高さ}
\newcommand\rxnman[2]{\begin{scope}[shift={#1}, scale=#2/10]
  \filldraw[fill=bkYellow!25, draw=bkInk, line width=0.5pt] (0,8.6) circle (1.4);
  \draw[bkInk, line width=0.8pt] (0,7.2) -- (0,3.4) (0,3.4) -- (-1.3,0) (0,3.4) -- (1.3,0) (-1.8,5.6) -- (0,6.6) -- (1.8,5.6);
\end{scope}}
% 目（右向き）  \rxneye{(x,y)}
\newcommand\rxneye[1]{\begin{scope}[shift={#1}]
  \draw[bkInk, line width=0.6pt] (-3,0) .. controls (-1,2.2) and (1,2.2) .. (2.6,0) .. controls (1,-2.2) and (-1,-2.2) .. (-3,0);
  \fill[bkInk] (0.2,0) circle (0.9);
\end{scope}}
% 偏光板（縦長の板と、通す向きの線）  \rxnpol{(x,y)}{通す向きの角（度、たてを0）}
\newcommand\rxnpol[2]{\begin{scope}[shift={#1}]
  \filldraw[fill=bkGray!18, draw=bkInk, line width=0.6pt] (-1.5,-9) -- (1.5,-7) -- (1.5,9) -- (-1.5,7) -- cycle;
  \begin{scope}[rotate=#2]\draw[bkInk, line width=0.9pt, {Stealth[length=2.4pt]}-{Stealth[length=2.4pt]}] (0,-5) -- (0,5);\end{scope}
\end{scope}}
\makeatother
\tikzset{rxn lab/.style={font=\sffamily\footnotesize, inner sep=1pt},
  rxn ray/.style={-{Stealth[length=3.4pt, width=3pt]}, line width=0.9pt, bkRed!80!black},
  rxn ray2/.style={-{Stealth[length=3.4pt, width=3pt]}, line width=0.9pt, bkBlue!80!black},
  rxn ray3/.style={-{Stealth[length=3.4pt, width=3pt]}, line width=0.9pt, bkInk},
  rxn dim/.style={{Stealth[length=2.6pt]}-{Stealth[length=2.6pt]}, line width=0.5pt, bkGray!80!black},
  rxn glass/.style={fill=bkBlue!10, draw=bkInk, line width=0.6pt}}

\phchapter{13}{光の波：\\干渉・回折・偏光・散乱}{空は、なぜ青い？　夕焼けは、なぜ赤い？}{\bklevelrange{3}{5}}{光の波：干渉・回折・偏光・散乱}
\bkmichishirube{4-9,5-17,7-10,9-4,熱-13,波-2,波-8,波-12}{波-14,電-15,原-2,原-4}

\begin{goal}
\begin{itemize}[label={\color{bkTheme}●}, leftmargin=1.6\zw, itemsep=0pt, topsep=1pt]
\item \bklevel{3}\enspace ヤングの実験・回折格子・薄い膜で、光が強め合う所と弱め合う所を、経路の差から計算できる。
\item \bklevel{5}\enspace 円い穴の回折で見分けられる角の限界と、偏光板を通る光の強さ・ブルースター角を計算できる。
\item \bklevel{5}\enspace ゆすられる電子のモデル（強制振動）から、空が青く夕焼けが赤いわけを説明できる。
\end{itemize}
\end{goal}

\phkoushiki{%
\phk{ヤングの実験 $\Delta x=L\lambda/d$}{高校物理}{ph:波-13:young}
\phk{回折格子 $d\sin\theta=m\lambda$}{高校物理}{ph:波-13:koushi}
\phk{薄い膜の干渉 $2nt\cos r=(m+\tfrac12)\lambda$ など}{高校物理}{ph:波-13:hakumaku}
\phk{くさび形の空気の層・ニュートンリング}{高校物理}{ph:波-13:kusabi}
\phk{偏光（光は横波）}{高校物理}{ph:波-13:henkou}
}
%% 整合（K-8・U-5）：「光の散乱（空が青い）」の行は公式ではなく、学習指導要領解説（平成30年告示）の「物理」(2)(ウ) の
%%   扱う事項（光の速さ・波長・反射・屈折・分散・偏光、ヤングの実験・回折格子・薄膜の干渉）にもないので、けずった（本文⑧はそのまま）。

\begin{hatena}
\noindent\begin{minipage}[c]{0.58\linewidth}\raggedright
晴れた日の空は青く、夕方の空は赤やだいだい色に染まります。でも、太陽の光そのものは、白い光です。空気は無色で、すき通っています。白い光と無色の空気から、なぜ昼は青、夕方は赤という、違う色が生まれるのでしょう。
\end{minipage}\hfill
\begin{minipage}[c]{0.39\linewidth}\centering
\begin{tikzpicture}[x=0.66mm, y=0.66mm]
  \fill[bkBlue!35] (0,0) rectangle (34,34);
  \rxnsun{(22,25)}{3.2}
  \fill[bkGreen!45] (0,0) rectangle (34,4);
  \rxnman{(10,4)}{9}
  \node[rxn lab, anchor=north] at (17,-1.5) {昼};
  \shade[top color=bkOrange!25, bottom color=bkRed!55] (40,0) rectangle (74,34);
  \begin{scope}\clip (40,4) rectangle (74,34); \rxnsun{(64,6)}{3.2}\end{scope}
  \fill[bkGreen!45!black!60] (40,0) rectangle (74,4);
  \rxnman{(50,4)}{9}
  \node[rxn lab, anchor=north] at (57,-1.5) {夕方};
\end{tikzpicture}
\end{minipage}
\end{hatena}

\noindent 光は波です。波の章（波-5〜8）で学んだ「重なる・回りこむ・はね返る」は、光でも起きます。ただ、光の波長は、髪の毛の太さの百分の1ほどしかありません。そのため、ふだんは波らしいふるまいが目に見えません。この章では、波長がとても短い光の「波らしさ」を、目で見える現象にしてとらえます。シャボン玉の色、CD の虹色、カメラのレンズの色、偏光サングラス、そして空の青。どれも、光が波であることから説明できます。

\noindent この章は3段に分かれています。①〜④（目でわかる・手でわかる）は、経路の差を数えるだけで、明るい所と暗い所が計算できます。⑤〜⑧（式でわかる）では、円い穴を通る光、光のゆれの向き、光を長さのものさしに使う方法、そして空の色を、波の式とゆれのモデルで調べます。

%% ---------------- やってみよう・思い出そう ----------------
\begin{yattemiyou}[コップの中に、夕焼けを作る]
\textbf{用意するもの}\par
透明なコップ（またはペットボトル）、水、牛乳（数てき）、スマホのライト、CD か DVD、偏光サングラス2つ（なければ1つと、パソコンやスマホの液晶の画面）
\begin{enumerate}[label={\bkcircled[bkGreen]{\arabic*}}, leftmargin=2.4\zw, itemsep=1pt, topsep=2pt]
\item コップに水を入れ、牛乳を2〜3てきたらしてかき混ぜる。暗い部屋で、ライトの光をコップの横から当てる。光の通り道を、横から見ると何色かな？
\item 光が来る向きの反対側から、コップごしにライトを見る（ライトを直接長く見つめない）。何色に見えるかな？　牛乳を増やすと、色はどう変わるかな？
\item CD の記録面に部屋の明かりを映し、少し傾けて、虹のような色の帯を見る。
\item 偏光サングラスを2つ重ね、片方だけを回しながら、明るい窓を見る。1つしかなければ、液晶の画面の前で回す。何度回すと、いちばん暗くなるかな？
\end{enumerate}
\end{yattemiyou}

\noindent ①②の結果は①と⑧で、③は③で、④は⑥で説明します。

\begin{omoidasou}
\bkrecall{波-5}{波の速さ＝振動数×波長。光では $c=f\lambda$。}
\bkrecall{熱-13}{光は電磁波。目に見える光（可視光）の波長は約 \qtyrange{0.38}{0.78}{\micro m}。熱-13 では比熱とぶつけないため光の速さを $c_0$ と書いたが、光の部のこの章では $c$ と書く。}
\bkrecall{波-8}{2つの波源の干渉：経路の差が波長の整数倍で強め合い、整数＋$\tfrac12$ 倍で弱め合う（波-8 ③）。すきまのはば $a$ の回折：はじめに暗くなる向きは $\sin\theta=\lambda/a$（波-8 ④）。}
\bkrecall{波-12}{屈折率 $n=c/v$。物質の中では光が $n$ 分の1の速さになる（波-12 ③）。}
\bkrecall{波-2}{ゆすり続けるときのゆれの振幅は、ゆする振動数が固有振動数から遠いと小さい（波-2 ④⑤）。}
\bkrecall{4-9}{同じ周期の波をたすと、同じ周期の1つの波になる（合成、4-9⑥）。}
\bkrecall{5-17}{小さな角 $\theta$（rad）では、$\sin\theta\approx\tan\theta\approx\theta$（5-17⑥）。}
\end{omoidasou}

%% ---------------- 物理の道具 ----------------
\bksection{doguu}{物理の道具}

\phstage{me}{光の波らしさを目で見る}{3}

\bksubsection{1}{光の色は波長で決まる：シャボン玉・CD・空}

\leavevmode\label{ph:波-13:sec1}熱-13 ①では、光が電気と磁気のゆれの波（電磁波）で、波長ごとに分けられることを学びました。そこでは「光の波長と色、光が波であることの詳しい話は『波動』で」と約束しました。この節から、その約束を果たします。

\Needspace{5\baselineskip}\noindent\textbf{波長と色}\par\nobreak
太陽の光をプリズムに通すと、紫・青・緑・黄・だいだい・赤の帯に分かれます（波-12 ⑧）。この色の帯は、波長の順に並んでいます。目に見える光（\textbf{可視光}）の波長は、約 \qty{380}{nm}（紫）から約 \qty{780}{nm}（赤）までです。\qty{1}{nm}（ナノメートル）は $10^{-9}$\,m、つまり \qty{1}{mm} の百万分の1です（波長と色の対応は図\ref{ph:波-13:fig-iro}）。

光の速さ $c$ は、真空の中で $c=\qty{299792458}{m/s}$（約 $3.00\times10^8$\,m/s）と決められています。$c=f\lambda$（波-5 ②）から、緑の光（\qty{550}{nm}）の振動数は
\pheq{f=\frac{c}{\lambda}=\frac{3.00\times10^8}{550\times10^{-9}}\approx5.45\times10^{14}\ \mathrm{Hz}}{Hz ＝ (m/s) ÷ m}
です。1秒に約545兆回ゆれています。音（数百 Hz）とは、まるでけたが違います。

\begin{phmijika}[光の波長は、どれくらい小さい？]
髪の毛の太さは約 \qty{0.08}{mm} です。緑の光の波長 \qty{550}{nm}＝\qty{0.00055}{mm} の波を並べると、髪の毛の太さの中に約145個も入ります。すきまや物の大きさが波長と同じくらいのときに、波はよく回りこみます（波-8 ④）。窓やドアの大きさは光の波長の百万倍以上あるので、光はほとんど回りこまず、まっすぐ進むように見えます。ドアの向こうの声は聞こえても、姿は見えないのはこのためです。
\end{phmijika}

\begin{bkfigure}
\begin{tikzpicture}[x=1mm, y=1mm]
  % x = (λ-380)*0.3
  \fill[bkPurple!55] (0,0) rectangle (21,9);
  \fill[bkBlue!60] (21,0) rectangle (34.5,9);
  \fill[bkGreen!60] (34.5,0) rectangle (57,9);
  \fill[bkYellow!85] (57,0) rectangle (63,9);
  \fill[bkOrange!75] (63,0) rectangle (72,9);
  \fill[bkRed!65] (72,0) rectangle (120,9);
  \draw[bkInk, line width=0.6pt] (0,0) rectangle (120,9);
  \foreach \l in {400,450,500,550,600,650,700,750}{\pgfmathsetmacro\xx{(\l-380)*0.3}\draw[bkInk] (\xx,0) -- (\xx,-1.5); \node[rxn lab, anchor=north] at (\xx,-2) {\l};}
  \node[rxn lab, anchor=north east] at (120,-6.5) {波長（nm）};
  \node[rxn lab, text=white] at (10.5,4.5) {紫};
  \node[rxn lab, text=white] at (27.75,4.5) {青};
  \node[rxn lab] at (45.75,4.5) {緑};
  \draw[bkInk, line width=0.4pt] (60,9) -- (60,13); \node[rxn lab, anchor=south] at (60,13.4) {黄};
  \draw[bkInk, line width=0.4pt] (67.5,9) -- (70,13); \node[rxn lab, anchor=south west] at (69,13.4) {だいだい};
  \node[rxn lab] at (96,4.5) {赤};
  \node[rxn lab, anchor=south west] at (0,10) {←紫外線};
  \node[rxn lab, anchor=south east] at (120,10) {赤外線→};
\end{tikzpicture}
\bkfigcaption{可視光の波長と色。色の境目は目安で、色は少しずつ移り変わっている。紫より短いと紫外線、赤より長いと赤外線（熱-13 ①）。}
\label{ph:波-13:fig-iro}
\end{bkfigure}

\Needspace{5\baselineskip}\noindent\textbf{光が波だと、色が生まれる}\par\nobreak
光の波長がこんなに小さくても、波らしさが目で見える場面があります。図\ref{ph:波-13:fig-sanshu}の3つです。

\begin{bkfigure}
\begin{tikzpicture}[x=1mm, y=1mm]
  % シャボン玉
  \begin{scope}[shift={(18,14)}]
    \shade[inner color=white, outer color=bkPurple!25] (0,0) circle (11);
    \draw[bkPurple!70, line width=1.1pt] (-9.5,4) arc (155:205:10.5);
    \draw[bkGreen!70!black, line width=1.1pt] (-7,7.5) arc (130:165:10);
    \draw[bkOrange, line width=1.1pt] (5,8) arc (60:30:10);
    \draw[bkBlue!80, line width=1.1pt] (8.5,-3) arc (-20:-55:10);
    \draw[bkInk, line width=0.5pt] (0,0) circle (11);
    \fill[white] (-4,5) ellipse (2 and 1.2);
  \end{scope}
  \node[rxn lab, align=center, anchor=north] at (18,0) {シャボン玉\\薄い膜の干渉（④）};
  % CD
  \begin{scope}[shift={(62,14)}]
    \filldraw[fill=bkGray!15, draw=bkInk, line width=0.6pt] (0,0) circle (12);
    \foreach \r/\c in {10.5/bkRed, 9.6/bkOrange, 8.7/bkYellow!90!black, 7.8/bkGreen, 6.9/bkBlue, 6.0/bkPurple}{\draw[\c, line width=1.3pt] (20:\r) arc (20:70:\r);}
    \filldraw[fill=white, draw=bkInk, line width=0.5pt] (0,0) circle (2.2);
  \end{scope}
  \node[rxn lab, align=center, anchor=north] at (62,0) {CD の虹色\\細かいみぞの回折（③）};
  % 空
  \begin{scope}[shift={(96,2)}]
    \fill[bkBlue!30] (0,0) rectangle (28,26);
    \rxnsun{(22,20)}{3}
    \foreach \a in {0,60,...,300}{\draw[bkBlue!80!black, line width=0.7pt, -{Stealth[length=2.2pt]}] (9,11) -- ++(\a:4.5);}
    \fill[bkInk] (9,11) circle (0.8);
    \draw[bkInk!60, line width=0.7pt, dashed] (18.5,18) -- (10.5,12);
  \end{scope}
  \node[rxn lab, align=center, anchor=north] at (110,0) {空の青\\空気の分子の散乱（⑧）};
\end{tikzpicture}
\bkfigcaption{光が波であることが、色として目に見える3つの場面。シャボン玉は膜の表と裏で反射した光の重なり、CD は細かいみぞで回りこんだ光の重なり、空の青は空気の分子が光を四方へまき散らすこと（散乱）で生まれる。}
\label{ph:波-13:fig-sanshu}
\end{bkfigure}

シャボン玉の色は、膜の表と裏ではね返った2つの光の\textbf{干渉}でできます（④）。CD の虹色は、細かいみぞに当たって回りこんだ光（\textbf{回折}）どうしの干渉です（③）。どちらも、「波長によって、強め合う向きや場所が違う」ので、白い光が色に分かれて見えます。3つめの空の色は、少し違うしくみです。

\Needspace{5\baselineskip}\noindent\textbf{空気の分子は、光をまき散らす}\par\nobreak
\label{ph:波-13:hatena1}太陽の光が空気の中を通るとき、ほとんどの光はまっすぐ進みます。でも、ほんの一部は、空気の分子（窒素や酸素）に当たって、四方八方へまき散らされます。これを光の\textbf{散乱}といいます。\bkidx{さんらん}{散乱（光の）}空を見上げたとき、太陽のない方向から光が来るのは、この散乱した光を見ているからです。月には空気がないので、昼でも空は真っ暗です。

1871年、イギリスのレイリーは、空気の分子のように波長よりずっと小さな粒が光を散乱する強さは、\textbf{波長の4乗に反比例する}ことを示しました。
\[ \text{散乱の強さ}\propto\frac{1}{\lambda^4} \]
\label{ph:波-13:rayleigh}これを\textbf{レイリー散乱}といいます。\bkidx{れいりーさんらん}{レイリー散乱}$\propto$ は「比例する」の記号です。なぜ4乗になるのかは、⑧で、ゆすられる電子のモデルから導きます。

\begin{phyomi}[散乱の強さ $\propto 1/\lambda^4$]
「波長が短い光ほど、ずっと強く散らされる」。青（\qty{450}{nm}）は赤（\qty{700}{nm}）より波長が約 0.64 倍なので、散らされる強さは $(700/450)^4\approx5.9$ 倍です（\nobreak\bkreidairef{1}\nobreak）。
\end{phyomi}

\begin{phbaisuru}{波長 $\lambda$}
散乱の強さは $1/2^4=1/16$ になります。\qty{400}{nm}（紫）と \qty{800}{nm}（赤外線のはじまり）では、16倍も違います。グラフは、\qty{700}{nm} の赤の散乱を1としたときの、ほかの色の散乱の強さです。
\par\vspace{2pt}\centering
\begin{tikzpicture}
\begin{axis}[bkplot, width=104mm, height=46mm, xmin=380, xmax=780, ymin=0, ymax=12, xtick={400,450,500,550,600,650,700,750}, ytick={0,2,4,6,8,10}, xlabel={$\lambda$（nm）}, ylabel={散乱の強さ（\qty{700}{nm} を1）}]
  \addplot[bkcurve, domain=380:780, samples=100] {(700/x)^4};
  \addplot[only marks, mark=*, mark size=1.8pt, bkOrange] coordinates {(450,5.86) (550,2.62) (700,1)};
  \node[font=\sffamily\footnotesize, anchor=south west] at (axis cs:452,5.9) {青 5.9};
  \node[font=\sffamily\footnotesize, anchor=south west] at (axis cs:552,2.7) {緑 2.6};
  \node[font=\sffamily\footnotesize, anchor=south west] at (axis cs:702,1.1) {赤 1};
\end{axis}
\end{tikzpicture}
\end{phbaisuru}

\Needspace{5\baselineskip}\noindent\textbf{はてな？の答え（目で見る）}\par\nobreak
\bkhatenaref の答えを、図\ref{ph:波-13:fig-yuuyake}で見ましょう。昼、太陽の光が空気の中を通るとき、青い光は赤い光の約6倍も強く、横へ散らされます。どの方向の空を見ても、目に入ってくるのは、空気の分子が散らした光です。それは青い光を多くふくむので、空は青く見えます。

夕方、太陽が地平線の近くにあると、光は空気の中を、昼よりずっと長い道のりで通ってきます。その途中で、青い光は次々に横へ散らされて、なくなってしまいます。目まで届くのは、散らされにくい赤やだいだいの光です。だから、夕日と、その周りの空は赤く見えるのです（数で確かめるのは\bkreidairef{13}）。

\begin{bkfigure}
\begin{tikzpicture}[x=1mm, y=1mm]
  % 地球と大気（たてに大きく引きのばしてある）
  \fill[bkBlue!12] (0,0) arc (180:0:60 and 34) -- (114,0) arc (0:180:54 and 28) -- cycle;
  \fill[bkGreen!35] (6,0) arc (180:0:54 and 28) -- cycle;
  \draw[bkInk, line width=0.5pt] (6,0) arc (180:0:54 and 28);
  \draw[bkBlue!60, line width=0.5pt, dashed] (0,0) arc (180:0:60 and 34);
  \node[rxn lab] at (60,12) {地球};
  \node[rxn lab, anchor=west] at (96,33) {空気の層};
  \draw[bkGray, line width=0.4pt] (95.5,32.5) -- (90,26.5);
  % 昼の人（てっぺん）
  \rxnman{(60,28)}{6}
  \rxnsun{(60,52)}{3}
  \draw[rxn ray, bkOrange!90!black] (60,47.5) -- (60,35);
  \node[rxn lab, anchor=west] at (61.5,44) {昼：短い道のり};
  % 夕方の人（右はし）
  \rxnman{(21,19.5)}{6}
  \rxnsun{(-6,40)}{3}
  \draw[rxn ray] (-1.5,37.5) -- (18.5,24.5);
  \node[rxn lab, anchor=south west, align=left] at (-8,47) {夕方：長い道のり};
  % 散らされる青
  \foreach \p in {(4,33.9),(9,30.7),(14,27.5)}{\draw[bkBlue!80!black, line width=0.7pt, -{Stealth[length=2.2pt]}] \p -- ++(80:4.5);}
  \node[rxn lab, text=bkBlue!80!black, anchor=west] at (18,36) {青が散らされて減る};
\end{tikzpicture}
\bkfigcaption{昼と夕方の、太陽の光が通る空気の道のり（空気の層の厚さは、たてに大きく引きのばしてある）。夕方は道のりが長く、そのあいだに青い光が散らされてしまい、赤い光が残る。}
\label{ph:波-13:fig-yuuyake}
\end{bkfigure}

やってみようの①②で、牛乳を少し入れた水を横から見ると青っぽく、正面からライトを見るとだいだい色に見えたのは、同じしくみです。牛乳の小さな粒が光を散乱し、横へは青い光が多く出て、まっすぐ通りぬける光は赤っぽくなります。

\begin{seikaku}
牛乳の脂肪の粒は、空気の分子よりずっと大きく、光の波長と同じくらいのものも多くふくみます。そのため、散乱の強さは $1/\lambda^4$ にはぴったり従わず、色の分かれ方は空よりおだやかです。それでも、小さな粒ほど青い光をよく散らすという傾向は同じです。雲が白いのは、雲の水の粒（直径 数〜数十 \unit{\micro m}）が波長よりずっと大きく、どの色もほぼ同じように散らすからです（⑧）。
\end{seikaku}

\begin{motto}[3]
\phboxhead{火星の夕焼けは、青い}
火星の探査車が写した写真では、火星の昼の空は薄い赤茶色で、夕方には、しずむ太陽の周りだけが青白く見えます。火星の大気はとても薄いので、空気の分子による散乱は弱く、空の色は、大気に浮かぶ細かいちり（直径 \unit{\micro m} ほど）で決まります。このちりは、青い光を吸いやすく、赤っぽい光を広く散らすので、昼の空は赤茶色になります。一方、青い光は、太陽の進む向きに近い向きへ散らされやすいので、夕方、太陽の方向を見ると、その周りが青く見えると考えられています。空の色は、「何が光を散らすか」で決まるのです。
\end{motto}

\Needspace{12\baselineskip}
\phstage{te}{経路の差で明るさを計算する}{3}

\bksubsection{2}{ヤングの実験：光が波である証拠}

\leavevmode\label{ph:波-13:sec2}光は、粒の流れなのでしょうか、それとも波なのでしょうか。17世紀の終わりごろ、ニュートンは光を粒の流れと考え、ホイヘンスは波と考えていました（波-8 ①のホイヘンスの原理）。決め手になったのは、1801年から1803年ごろに、イギリスのトーマス・ヤングが講演と論文で示した実験です。ヤングは、力-8 で学んだヤング率のヤングと同じ人です。

\Needspace{5\baselineskip}\noindent\textbf{実験のしかた}\par\nobreak
図\ref{ph:波-13:fig-young}のように、1つの光源からの光を、間かく $d$ の2つの細いすきま（スリット）に当てます。2つのすきまを通った光が重なる所に、距離 $L$ だけ離してスクリーンを置きます。すると、スクリーンには、明るい線と暗い線が、同じ間かくで交互に並びます。これを\textbf{干渉じま}といいます。\bkidx{かんしょうじま}{干渉じま}\bkidx{やんぐのじっけん}{ヤングの実験}

\begin{bkfigure}
\begin{tikzpicture}[x=1mm, y=1mm]
  % レーザー
  \filldraw[fill=bkGray!35, draw=bkInk, line width=0.6pt] (0,11) rectangle (16,17);
  \node[rxn lab] at (8,14) {レーザー};
  \draw[bkRed!80!black, line width=1.6pt] (16,14) -- (33,14);
  % すきまの板
  \fill[bkInk!80] (33,2) rectangle (35,12.6);
  \fill[bkInk!80] (33,13.4) rectangle (35,14.6);
  \fill[bkInk!80] (33,15.4) rectangle (35,26);
  % 広がる光
  \foreach \e in {-11,-6,0,6,11}{\draw[bkRed!60, line width=0.5pt] (35,13) -- (92,14+\e); \draw[bkRed!60, line width=0.5pt] (35,15) -- (92,14+\e);}
  % スクリーン（横から）
  \fill[bkGray!40] (92,0) rectangle (94,28);
  \draw[rxn dim] (34,-2) -- (93,-2);
  \node[rxn lab, anchor=north] at (63.5,-2.8) {$L$（約 \qty{1}{m}）};
  \draw[bkInk, line width=0.4pt] (34,13) -- (27,22);
  \node[rxn lab, anchor=south east, align=right] at (28,22) {2つのすきま\\（間かく $d$）};
  \node[rxn lab, anchor=south] at (93,28.6) {スクリーン};
  % スクリーン（正面から）
  \begin{scope}[shift={(104,0)}]
    \fill[bkInk!85] (0,0) rectangle (22,28);
    \foreach \k in {-4,...,4}{\shade[top color=bkInk!85, bottom color=bkInk!85, middle color=bkRed!75] (0,14+3.2*\k-1.1) rectangle (22,14+3.2*\k+1.1);}
    \draw[bkInk, line width=0.6pt] (0,0) rectangle (22,28);
    \draw[rxn dim, bkInk] (24,14) -- (24,17.2);
    \node[rxn lab, anchor=west] at (25,15.6) {$\Delta x$};
    \node[rxn lab, anchor=south] at (11,28.6) {正面から見る};
  \end{scope}
\end{tikzpicture}
\bkfigcaption{ヤングの実験。2つのすきまを通った光が重なって、スクリーンに明るい線と暗い線が同じ間かくで並ぶ（右は、スクリーンを正面から見たところ）。すきまの間かくは \qty{0.2}{mm} ほどで、図では大きくかいてある。}
\label{ph:波-13:fig-young}
\end{bkfigure}

もし光が粒の流れなら、2つのすきまの正面だけが明るく光るはずです。ところが実際には、すきまの正面から外れた所にも明るい線が何本もでき、そのあいだに暗い線ができます。「光と光を重ねると、暗くなる所がある」。これは、山と谷が打ち消し合う波でしか起きません（波-8 ③）。こうして、光が波であることが確かめられました。

\Needspace{5\baselineskip}\noindent\textbf{経路の差を数える}\par\nobreak
スクリーンの上の点 P の明るさは、2つのすきまから P までの距離の差（\textbf{経路の差}）で決まります。波-8 ③と同じように、経路の差が波長の整数倍なら山と山が重なって強め合い、整数＋$\tfrac12$ 倍なら山と谷が重なって弱め合います。

\begin{bkfigure}
\begin{tikzpicture}[x=1mm, y=1mm]
  \coordinate (S1) at (10,24); \coordinate (S2) at (10,14); \coordinate (M) at (10,19);
  \coordinate (P) at (118,44); \coordinate (O) at (118,19); \coordinate (H) at (12.58,14.72);
  \fill[bkGray!40] (118,6) rectangle (120,50);
  \fill[bkInk!80] (8,4) rectangle (10,13.4); \fill[bkInk!80] (8,14.6) rectangle (10,23.4); \fill[bkInk!80] (8,24.6) rectangle (10,34);
  \draw[bkRed!80!black, line width=0.8pt] (S1) -- (P);
  \draw[bkRed!80!black, line width=0.8pt] (S2) -- (P);
  \draw[bkGray, line width=0.5pt, dashed] (M) -- (O);
  \draw[bkGray, line width=0.5pt, dashed] (M) -- (P);
  \draw[bkInk, line width=0.6pt] (S1) -- (H);
  \draw[bkBlue!80!black, line width=1.6pt] (S2) -- (H);
  \draw[bkInk, line width=0.4pt] (12.15,16.26) -- ++(1.54,0.43) -- ++(0.43,-1.54);
  \fill[bkInk] (S1) circle (0.6); \fill[bkInk] (S2) circle (0.6); \fill[bkInk] (P) circle (0.8); \fill[bkInk] (O) circle (0.8);
  \node[rxn lab, anchor=east] at (6.3,25.5) {S$_1$};
  \node[rxn lab, anchor=east] at (6.3,12.5) {S$_2$};
  \draw[rxn dim] (1.5,14) -- (1.5,24); \node[rxn lab, anchor=east] at (1,19) {$d$};
  \node[rxn lab, anchor=west] at (122,44) {P};
  \node[rxn lab, anchor=west] at (122,19) {O};
  \draw[rxn dim] (125,19.5) -- (125,43.5); \node[rxn lab, anchor=west] at (125.6,31.5) {$x$};
  \draw[rxn dim] (10,1.5) -- (118,1.5); \node[rxn lab, fill=white, inner sep=1.5pt] at (64,1.5) {$L$};
  \draw[bkInk, line width=0.5pt] (36,19) arc (0:13:26); \node[rxn lab] at (42.5,21.8) {$\theta$};
  \draw[bkBlue!80!black, line width=0.4pt] (11.6,14.2) -- (22,8);
  \node[rxn lab, text=bkBlue!80!black, anchor=west, align=left] at (22.5,8) {経路の差 $d\sin\theta$};
\end{tikzpicture}
\bkfigcaption{経路の差の作り方。S$_1$ から S$_2$P の線へ垂線を下ろすと、S$_2$P のほうが太い線（$d\sin\theta$）だけ長い。実際は $L$ が $d$ の約5000倍あるので、S$_1$P と S$_2$P はほぼ平行になる。図は、たてに大きく引きのばしてある。}
\label{ph:波-13:fig-keiro}
\end{bkfigure}

図\ref{ph:波-13:fig-keiro}で、スクリーンはすきまからとても遠い（$L\gg d$）ので、S$_1$P と S$_2$P はほぼ平行です。S$_1$ から S$_2$P へ垂線を下ろすと、2つの経路の差は、太い線の長さ
\pheq{\text{経路の差}\approx d\sin\theta}{m ＝ m}
です。ここで $\theta$ は、すきまのまん中から見た P の向きの角です（直角三角形の角の関係で、S$_1$ の所の角も $\theta$ になる）。

P の位置を、スクリーンのまん中 O からの距離 $x$ で表すと、$\tan\theta=x/L$ です。$\theta$ はとても小さい（\qty{0.01}{rad} ほど）ので、$\sin\theta\approx\tan\theta$（5-17⑥）を使って
\[ \text{経路の差}\approx d\sin\theta\approx d\tan\theta=\frac{dx}{L} \]
となります。明るい線ができるのは、これが波長の整数倍になる所です。
\[ \frac{dx}{L}=m\lambda\qquad(m=0,\ \pm1,\ \pm2,\ \cdots) \]
$m$ を\textbf{次数}といいます。\bkidx{じすう}{次数}$x$ について解くと、$m$ 番めの明るい線の位置は $x_m=mL\lambda/d$ です。となりどうしの明るい線の間かくは、$x_{m+1}-x_m$ なので\label{ph:波-13:young}%% ページの数を式の始まりに合わせる（式の前では改ページしない。修3-6、P6-17）
\pheq{\Delta x=\frac{L\lambda}{d}}{m ＝ m × m ÷ m}
となります。暗い線は、明るい線と明るい線のちょうどまん中（経路の差が $(m+\tfrac12)\lambda$ の所）にできます。

\begin{phyomi}[$\Delta x=L\lambda/d$]
「しまの間かくは、スクリーンが遠いほど広く、波長が長いほど広く、すきまの間かくが広いほどせまい」。$L=\qty{1}{m}$、$d=\qty{0.2}{mm}$ なら、$L/d=5000$ なので、波長 \qty{633}{nm} が5000倍に拡大されて、約 \qty{3.2}{mm} の間かくとして目に見えます（\nobreak\bkreidairef{2}\nobreak）。ヤングの実験は、見えないほど小さい波長を、目ではかれる長さに引きのばす「ものさし」なのです。
\end{phyomi}

\Needspace{10\baselineskip}
\begin{phflow}[ヤングの実験]
\phflowrow{2つのすきまからの\\光の重なり\\経路の差 $d\sin\theta=m\lambda$}{スクリーンが遠い $L\gg d$\\小さな角 $\sin\theta\approx x/L$}{$\Delta x=\dfrac{L\lambda}{d}$}
\end{phflow}

\begin{phmodel}[ヤングの実験]
\phignore{すきまのはば（細いとする）。2つのすきまを通る光の強さの違い。空気の屈折率（1.0003 を1とする）。}
\phassume{1つの波長の光。2つのすきまの光は、同じ光源から分かれたので、山と谷のそろい方（位相）がそろっている。スクリーンは十分遠い。}
\phlimit{すきまのはばが広いと、1つのすきまの中の回折（⑤）で、明るい線の明るさが外側ほど弱まる。スクリーンが近いと、経路の差を $d\sin\theta$ で表せない。白い光では、しまが色づく。}
\end{phmodel}

\Needspace{5\baselineskip}\noindent\textbf{表計算：明るさを並べる}\par\nobreak
明るい線とその間の明るさを、表計算で調べます。2つのすきまの波を $\sin\omega t$ と、経路の差のぶんおくれた $\sin(\omega t-\delta)$ とします（$\omega$ は光の角振動数、$t$ は時刻。④の膜の厚さ $t$ とは別の量です）。$\delta=2\pi\times(\text{経路の差}/\lambda)$ です（1波長おくれると $2\pi$）。積と和の式（4-9⑦）でたすと
\[ \sin\omega t+\sin(\omega t-\delta)=2\cos\frac{\delta}{2}\,\sin\Bigl(\omega t-\frac{\delta}{2}\Bigr) \]
となり、重なった波の振幅は $2\cos(\delta/2)$ です。光の明るさ（強さ）は振幅の2乗に比例する（波-7 ①）ので、まん中の明るさを1にすると、明るさは $\cos^2(\delta/2)=\cos^2(\pi\times\text{経路の差}/\lambda)$ になります。

\begin{bkcenter}
\bksheet{4}{26}{%
1 & $x$（mm） & 経路の差（nm） & 経路の差÷$\lambda$ & 明るさ \\
2 & 0 & 0 & 0.000 & 1.000 \\
3 & 0.8 & 160 & 0.253 & 0.491 \\
4 & 1.6 & 320 & 0.506 & 0.000 \\
5 & 2.4 & 480 & 0.758 & 0.526 \\
6 & 3.2 & 640 & 1.011 & 0.999 \\
7 & 4.0 & 800 & 1.264 & 0.457 \\
8 & 4.8 & 960 & 1.517 & 0.003 \\}
\bktabcaption{$d=\qty{0.20}{mm}$、$L=\qty{1.0}{m}$、$\lambda=\qty{633}{nm}$。B2 は \texttt{=0.20*A2/1000*1000000}（$dx/L$ を mm で出して nm に直す）、C2 は \texttt{=B2/633}、D2 は \texttt{=COS(PI()*C2)\^{}2}。}
\label{ph:波-13:tab-young}
\end{bkcenter}

経路の差÷$\lambda$ が 0、1 の所（$x=0$、約 \qty{3.2}{mm}）で明るさがほぼ1、0.5、1.5 の所（約 1.6、\qty{4.8}{mm}）でほぼ0になりました。明るい線の間かくは約 \qty{3.2}{mm} で、$L\lambda/d=\qty{3.165}{mm}$ と合います。

%% 短い囲みがまたがないように（修3-6、P6-11）
\Needspace{17\baselineskip}
\begin{phbaisuru}{すきまの間かく $d$}
しまの間かく $\Delta x$ は半分になります。グラフの実線は $d=\qty{0.20}{mm}$（$\Delta x=\qty{3.2}{mm}$）、破線は $d=\qty{0.40}{mm}$（$\Delta x=\qty{1.6}{mm}$）の明るさです。$d$ を広げると、しまがつまって見分けにくくなります。だから、ヤングの実験では、すきまをとても近づけて作ります。
\par\vspace{2pt}\centering
\begin{tikzpicture}
\begin{axis}[bkplot, axis lines=left, width=110mm, height=40mm, xmin=-8, xmax=8, ymin=0, ymax=1.25, xtick={-8,-6,-4,-2,0,2,4,6,8}, ytick={0,0.5,1}, xlabel={$x$（mm）}, ylabel={明るさ}]
  \addplot[bkcurve, domain=-8:8, samples=241] {cos(180*x/3.165)^2};
  \addplot[bkcurve2, dashed, domain=-8:8, samples=241] {cos(180*x/1.5825)^2};
\end{axis}
\end{tikzpicture}
\end{phbaisuru}

\begin{phmijika}[白い光で見るしま]
白い光でヤングの実験をすると、まん中（$m=0$）の明るい線は、どの波長でも経路の差が0なので白く見えます。そのほかの明るい線は、波長の短い青が内側、長い赤が外側にずれて、虹のように色づきます（$\Delta x$ が $\lambda$ に比例するから）。外へいくほど色が重なって、しまは見えなくなります。
\end{phmijika}

\begin{seikaku}
2つの電球を並べて照らしても、干渉じまはできません。電球の光は、たくさんの原子がばらばらに出した、とても短い波の列（長さは \unit{\micro m} 程度）の集まりで、2つの電球の光の山と谷のずれ方（位相の差）は、1秒に何兆回も入れかわってしまうからです。しまは一瞬ごとに動き回り、目には平均の明るさしか見えません。ヤングは、1つの光源の光を2つに分けることで、位相の差が変わらない（そろった）2つの光を作りました。レーザーの光は、波の列がとても長くそろっているので、干渉の実験がかんたんにできます。
\end{seikaku}

\bksubsection{3}{回折格子と分光：波長をはかる}

\leavevmode\label{ph:波-13:sec3}すきまを2つから、同じ間かく $d$ でたくさん（$N$ 本）に増やします。ガラスの板に、1 mm あたり何百本、何千本もの細い溝を刻んだものを\textbf{回折格子}といいます。\bkidx{かいせつこうし}{回折格子}光は、溝の間のすきま（または、溝と溝の間の平らな所）から回折して広がり、重なり合います。

\begin{bkfigure}
\begin{tikzpicture}[x=1mm, y=1mm]
  \fill[bkBlue!8] (18,0) rectangle (22,44);
  \foreach \k in {0,...,6}{\fill[bkInk!80] (18,1+6*\k) rectangle (22,1+6*\k+3.2);}
  \foreach \k in {0,...,6}{\draw[bkRed!80!black, line width=0.7pt, -{Stealth[length=2.6pt]}] (20,5+6*\k) -- ++(30:34);}
  \foreach \k in {0,...,6}{\draw[bkInk!60, line width=0.6pt, -{Stealth[length=2.4pt]}] (2,5+6*\k) -- (16.5,5+6*\k);}
  \node[rxn lab, anchor=south] at (8,44.5) {光};
  \draw[bkBlue!80!black, line width=1.6pt] (20,5) -- (22.6,6.5);
  \draw[bkInk, line width=0.5pt] (20,11) -- (22.6,6.5);
  \draw[bkInk, line width=0.4pt] (22.6,6.5) ++(150:1.4) -- ++(240:1.4) -- ++(330:1.4);
  \draw[bkBlue!80!black, line width=0.4pt] (21.7,5.6) -- (30,-1);
  \node[rxn lab, text=bkBlue!80!black, anchor=west] at (30.5,-1) {となりとの経路の差 $d\sin\theta$};
  \draw[rxn dim] (14.5,5) -- (14.5,11); \node[rxn lab, anchor=east] at (14,8) {$d$};
  \draw[bkGray, line width=0.5pt, dashed] (20,41) -- (50,41);
  \draw[bkInk, line width=0.5pt] (29,41) arc (0:30:9); \node[rxn lab, fill=white, inner sep=0.5pt] at (31,43.2) {$\theta$};
  \node[rxn lab, anchor=west, align=left] at (60,30) {どの2本の間でも、\\経路の差は同じ $d\sin\theta$。\\$d\sin\theta=m\lambda$ の向きでは、\\$N$ 本の光の山がそろい、\\強く明るくなる。};
\end{tikzpicture}
\bkfigcaption{回折格子。$N$ 本のすきまから、角 $\theta$ の向きへ進む光を考える（遠くで、またはレンズで集めて見る）。となりどうしの経路の差は、ヤングの実験と同じ $d\sin\theta$。}
\label{ph:波-13:fig-koushi}
\end{bkfigure}

図\ref{ph:波-13:fig-koushi}のように、角 $\theta$ の向きでは、となりどうしの経路の差がどれも $d\sin\theta$ です。これが波長の整数倍なら、1本め、2本め、…、$N$ 本めの光が、すべて山と山をそろえて重なります。だから、明るい線ができる向きは\label{ph:波-13:koushi}%
\pheq{d\sin\theta=m\lambda\qquad(m=0,\ \pm1,\ \pm2,\ \cdots)}{m ＝ m}
です。式の形はヤングの実験の明るい線と同じですが、違うのは、$N$ が大きいと、明るい線がとても細く、鋭くなることです。

\begin{phyomi}[$d\sin\theta=m\lambda$]
「となりの溝との経路の差が、ちょうど波長の $m$ 倍になる向きにだけ、光が強く出る」。$\sin\theta\le1$ なので、$m\lambda\le d$ の次数までしか現れません。波長が長い（赤い）光ほど、大きな角に曲がります。プリズム（波-12 ⑧）では赤が曲がりにくく紫がよく曲がるのと、反対の並びです。
\end{phyomi}

\Needspace{5\baselineskip}\noindent\textbf{矢印をつないで、重なりを見る}\par\nobreak
同じ周期の波をたすことは、回る矢印をつなぐことと同じです（4-9⑥の合成、6-10②）。矢印の長さが振幅、向きが山のずれ（位相）です。となりのすきまの光は、位相が $2\pi\times(d\sin\theta/\lambda)$ ずつずれるので、矢印は少しずつ向きを変えながらつながります。

\begin{bkfigure}
\begin{tikzpicture}[x=1mm, y=1mm]
  \tikzset{arw/.style={bkRed!80!black, line width=1.1pt, -{Stealth[length=2.6pt, width=2.4pt]}}, res/.style={bkBlue!80!black, line width=0.8pt, dashed, -{Stealth[length=2.8pt]}}}
  % ずれ0
  \begin{scope}[shift={(0,6)}]
    \foreach \k in {0,...,4}{\draw[arw] (7*\k,0) -- (7*\k+7,0);}
    \draw[res] (0,-3) -- (35,-3);
  \end{scope}
  \node[rxn lab, align=center, anchor=north] at (17.5,-4) {ずれ0\\合計の長さ 5（明るさ 1）};
  % ずれ 0.1 周
  \begin{scope}[shift={(55,0)}]
    \draw[arw] (0,0) -- (7,0); \draw[arw] (7,0) -- (12.66,4.11); \draw[arw] (12.66,4.11) -- (14.82,10.77);
    \draw[arw] (14.82,10.77) -- (12.66,17.43); \draw[arw] (12.66,17.43) -- (7.0,21.54);
    \draw[res] (0,0) -- (6.8,20.9);
  \end{scope}
  \node[rxn lab, align=center, anchor=north] at (62,-4) {ずれ 0.1 周\\合計の長さ 2.3（明るさ 0.42）};
  % ずれ 0.2 周
  \begin{scope}[shift={(102,3)}]
    \draw[arw] (0,0) -- (7,0); \draw[arw] (7,0) -- (9.16,6.66); \draw[arw] (9.16,6.66) -- (3.5,10.77);
    \draw[arw] (3.5,10.77) -- (-2.16,6.66); \draw[arw] (-2.16,6.66) -- (0,0);
    \fill[bkBlue!80!black] (0,0) circle (0.9);
  \end{scope}
  \node[rxn lab, align=center, anchor=north] at (105.5,-4) {ずれ 0.2 周\\もとにもどって0（暗い）};
\end{tikzpicture}
\bkfigcaption{5本のすきまの光を、矢印（長さ1を 7 mm でかいた）でつなぐ。となりとの位相のずれが0なら矢印はまっすぐ並んで長さ5、0.2 周（$\tfrac15$ 周）ずれると5本で1周してもとにもどり、合計は0になる。破線の矢印が合計。明るさは合計の長さの2乗を $5^2$ でわった値。}
\label{ph:波-13:fig-yajirushi}
\end{bkfigure}

図\ref{ph:波-13:fig-yajirushi}から、$N$ 本の矢印は、位相のずれが $\tfrac1N$ 周になると、ちょうど1周してもとにもどり、合計が0になることがわかります。$N$ が大きいほど、少しずれただけで暗くなります。これを Python で確かめます。矢印の和は、$x$ 成分と $y$ 成分に分けてたします（4-11）。

\begin{phcode}[Python]
from math import cos, sin, pi
def akarusa(N, u):        \# u：となりとの経路の差 ÷ 波長
    x = y = 0.0
    for p in range(N):    \# p 番めのすきまの矢印（長さ1）
        x = x + cos(2 * pi * p * u)
        y = y + sin(2 * pi * p * u)
    return (x * x + y * y) / N**2    \# まん中の明るさを1にする
us = [0, 0.05, 0.1, 0.25, 0.5, 1.0]
for N in [2, 5, 20]:
    print(N, [round(akarusa(N, u), 3) for u in us])
\end{phcode}
\begin{phcodeout}
2 [1.0, 0.976, 0.905, 0.5, 0.0, 1.0]
5 [1.0, 0.817, 0.419, 0.04, 0.04, 1.0]
20 [1.0, 0.0, 0.0, 0.0, 0.0, 1.0]
\end{phcodeout}

$u=d\sin\theta/\lambda$ が0と1（$m=0$ と $m=1$）では、何本でも明るさ1です。そのあいだは、$N=2$ ではゆるやかに変わります（ヤングの実験の $\cos^2$）が、$N=20$ では、$u=0.05$ でもう0です。図\ref{ph:波-13:fig-koushi-graph}のように、$N$ を増やすほど、明るい線ははばがせまく、鋭くなり、そのあいだはほぼまっ暗になります。

\begin{bkfigure}
\begin{tikzpicture}
\begin{axis}[bkplot, width=120mm, height=50mm, xmin=0, xmax=2.05, ymin=0, ymax=1.18, xtick={0,0.5,1,1.5,2}, ytick={0,0.5,1}, xlabel={$d\sin\theta/\lambda$}, ylabel={明るさ（山を1）}]
  \addplot[bkcurve3, densely dotted, line width=1.5pt, domain=0.0005:2.0495, samples=400] {(sin(180*2*x)/(2*sin(180*x)))^2};
  \addplot[bkcurve2, dashed, domain=0.0005:2.0495, samples=400] {(sin(180*5*x)/(5*sin(180*x)))^2};
  \addplot[bkcurve, domain=0.0005:2.0495, samples=800] {(sin(180*20*x)/(20*sin(180*x)))^2};
  \node[font=\sffamily\footnotesize, anchor=west, align=left, fill=white, inner sep=1.5pt] at (axis cs:0.30,0.95) {点線 $N=2$\\破線 $N=5$\\実線 $N=20$};
  \node[font=\sffamily\footnotesize, anchor=south] at (axis cs:1,1.03) {$m=1$};
  \node[font=\sffamily\footnotesize, anchor=south] at (axis cs:2,1.03) {$m=2$};
\end{axis}
\end{tikzpicture}
\bkfigcaption{すきまの数 $N$ と明るさ。横軸はとなりとの経路の差÷波長。整数の所（$d\sin\theta=m\lambda$）で明るく、$N$ が大きいほどその山は細い。山のはばは $N$ にほぼ反比例する。}
\label{ph:波-13:fig-koushi-graph}
\end{bkfigure}

\Needspace{5\baselineskip}\noindent\textbf{分光：光を波長ごとに分ける}\par\nobreak
白い光を回折格子に当てると、$d\sin\theta=m\lambda$ の $\theta$ は波長ごとに違うので、色ごとに別の向きへ進みます（図\ref{ph:波-13:fig-bunkou}）。$m=0$ はどの色も $\theta=0$ なので白、$m=\pm1$ では内側が紫、外側が赤の虹の帯になります。明るい線が細いので、波長が少し違う光も、別の向きにはっきり分かれます。こうして光を波長ごとに分けることを\textbf{分光}といい、そのための道具を\textbf{分光器}といいます。\bkidx{ぶんこう}{分光}

熱-13 ③の図（黒体の光の、波長ごとの強さ）は、こうした分光器ではかったものです。分光器で、星の光を波長ごとに分けると、星の温度や、星にふくまれる元素がわかります。元素は、それぞれ決まった波長の光だけを出したり吸ったりするからです（原-4 で学びます）。工場では、物質がどの波長の光を吸うかをはかって、成分を調べます（分光分析）。

\begin{phmijika}[CD・DVD の虹色]
CD の記録面には、データを記録した点の列（トラック）が、\qty{1.6}{\micro m} の間かくで、うずまき状に並んでいます。この細かい並びが、反射光の回折格子として働きます。やってみようの③で虹色が見えたのはこのためです。間かくが \qty{0.74}{\micro m} の DVD では、同じ色が CD より大きな角に曲がります（\nobreak\bkreidairef{4}\nobreak）。
\end{phmijika}

\begin{bkfigure}
\begin{tikzpicture}[x=1mm, y=1mm]
  \begin{scope}[xscale=1.25, xshift=-8mm]
  \fill[bkGray!40] (54,0) rectangle (66,2);
  \node[rxn lab, anchor=north] at (60,-0.5) {回折格子（1 mm に500本）};
  \draw[bkInk!60, line width=1.4pt, -{Stealth[length=3pt]}] (60,-11) -- (60,-5.5);
  \node[rxn lab, anchor=east] at (59,-9) {白い光};
  \draw[bkInk, line width=0.7pt] (5,40) -- (115,40);
  \fill[bkGray!20] (58.6,40) rectangle (61.4,42.5); \draw[bkInk, line width=0.4pt] (58.6,40) rectangle (61.4,42.5);
  \node[rxn lab, anchor=south] at (60,43) {$m=0$};
  % m=1：400 nm 8.2、700 nm 14.9（y=40）
  \foreach \s in {1,-1}{
    \shade[left color=bkPurple!70, right color=bkRed!70, shading angle=\s*90-90] ({60+\s*8.2},40) rectangle ({60+\s*14.9},42.5);
    \draw[bkPurple!70!black, line width=0.4pt] (60,2) -- ({60+\s*8.2},40);
    \draw[bkRed!70!black, line width=0.4pt] (60,2) -- ({60+\s*14.9},40);
    \shade[left color=bkPurple!70, right color=bkRed!70, shading angle=\s*90-90] ({60+\s*17.5},40) rectangle ({60+\s*39.2},42.5);
    \draw[bkPurple!70!black, line width=0.4pt] (60,2) -- ({60+\s*17.5},40);
    \draw[bkRed!70!black, line width=0.4pt] (60,2) -- ({60+\s*39.2},40);
  }
  \node[rxn lab, anchor=south] at (71.5,43) {$m=1$};
  \node[rxn lab, anchor=south] at (88.4,43) {$m=2$};
  \node[rxn lab, anchor=south] at (48.5,43) {$m=-1$};
  \node[rxn lab, anchor=south] at (31.6,43) {$m=-2$};
  \node[rxn lab, anchor=west, align=left] at (92,24) {紫が内側\\赤が外側};
  \end{scope}
\end{tikzpicture}
\bkfigcaption{回折格子で白い光を分ける（1 mm に500本。帯の位置は $d\sin\theta=m\lambda$ で計算した）。$m=0$ はどの色も重なって白い。$m=2$ の帯は $m=1$ より広く、さらに外の $m=3$ の帯（図にない）と一部が重なる。}
\label{ph:波-13:fig-bunkou}
\end{bkfigure}

\begin{motto}[5]
\phboxhead{$N$ 本の明るさの式を、等比級数で導く}
となりとの位相のずれを $\varphi=2\pi d\sin\theta/\lambda$ とします。$p$ 番めのすきまの光を、回る矢印 $e^{jp\varphi}$（6-10。$j$ は虚数単位）で表すと、$N$ 本の合計は、公比 $e^{j\varphi}$ の等比級数（7-10⑤）です。
\[ \sum_{p=0}^{N-1}e^{jp\varphi}=\frac{1-e^{jN\varphi}}{1-e^{j\varphi}} \]
$|1-e^{j\alpha}|=|e^{-j\alpha/2}-e^{j\alpha/2}|=2\lvert\sin(\alpha/2)\rvert$（6-10⑧）を分子と分母に使うと、合計の長さは $\lvert\sin(N\varphi/2)\rvert/\lvert\sin(\varphi/2)\rvert$ です。山（$\varphi=0$）の長さ $N$ でわって2乗すると
\[ \text{明るさ}=\Bigl(\frac{\sin(N\varphi/2)}{N\sin(\varphi/2)}\Bigr)^2 \]
となります。図\ref{ph:波-13:fig-koushi-graph}は、この式のグラフです。$N=2$ では、2倍角の式（4-9④）で $\cos^2(\varphi/2)$ になり、ヤングの実験の明るさにもどります。はじめて0になるのは $N\varphi/2=\pi$、つまり $d\sin\theta/\lambda=1/N$ の所です。
\end{motto}

\bksubsection{4}{薄い膜の干渉：シャボン玉・反射防止・ニュートンリング}

\leavevmode\label{ph:波-13:sec4}シャボン玉や、水たまりに浮いた油の膜は、虹のような色にかがやきます。膜の表面で反射した光と、膜の裏の面で反射した光が重なって、干渉するからです（図\ref{ph:波-13:fig-maku}）。ヤングの実験との大きな違いは2つあります。1つは、片方の光が膜の中を通ること、もう1つは、反射のときに山と谷が入れかわることがあることです。

\begin{bkfigure}
\begin{tikzpicture}[x=1mm, y=1mm]
  \fill[bkBlue!10] (0,10) rectangle (80,22);
  \draw[bkInk, line width=0.7pt] (0,22) -- (80,22); \draw[bkInk, line width=0.7pt] (0,10) -- (80,10);
  \node[rxn lab, anchor=west] at (82,28) {空気（$n=1$）};
  \node[rxn lab, anchor=west] at (82,16) {膜（屈折率 $n$）};
  \node[rxn lab, anchor=west] at (82,4) {空気};
  \draw[rxn ray3] (6,40) -- (19.6,22.4);
  \draw[rxn ray] (20.4,22.4) -- (34,40); \node[rxn lab, text=bkRed!80!black, anchor=south] at (34,40.5) {反射光1};
  \draw[bkBlue!70!black, line width=0.9pt] (20,22) -- (26,10.4);
  \draw[bkBlue!70!black, line width=0.9pt] (26.5,10.4) -- (32,22);
  \draw[rxn ray2] (32,22) -- (45.6,40); \node[rxn lab, text=bkBlue!80!black, anchor=south] at (47,40.5) {反射光2};
  \draw[rxn dim] (-2,10) -- (-2,22); \node[rxn lab, anchor=east] at (-2.5,16) {$t$};
  \node[rxn lab] at (26.6,24.8) {①};
  \node[rxn lab, anchor=north] at (26.3,8.6) {②};
  \node[rxn lab, anchor=north west, align=left] at (0,1.5) {①表の面：屈折率の大きい膜との境で反射→山と谷が入れかわる\\②裏の面：屈折率の小さい空気との境で反射→そのまま};
\end{tikzpicture}
\bkfigcaption{シャボン膜の2つの反射。反射光2は膜の中を往復するぶん（ほぼ $2t$）だけ長く進む。図ではわかりやすくするため、光をななめにかいた（ほぼ垂直に見るときの式を先に作る）。}
\label{ph:波-13:fig-maku}
\end{bkfigure}

\Needspace{5\baselineskip}\noindent\textbf{膜の中では、波長が $n$ 分の1になる}\par\nobreak
屈折率 $n$ の物質の中では、光の速さが $c/n$ になります（波-12 ③）。光が物質に入っても振動数 $f$ は変わらない（波-8 ①）ので、$v=f\lambda'$ から、物質の中の波長は $\lambda'=\lambda/n$ に縮みます。厚さ $t$ の膜を往復する長さ $2t$ の中には、$2t\div(\lambda/n)=2nt/\lambda$ 個の波が入ります。これは、真空の中を $2nt$ 進んだときと同じ数です。そこで、物質の中の長さに屈折率をかけた $n\times(\text{長さ})$ を\textbf{光路長}とよび、経路の差は光路長の差で数えます。\bkidx{こうろちょう}{光路長}膜を往復する光の、光路長の差は $2nt$ です。

\Needspace{5\baselineskip}\noindent\textbf{反射で、山と谷が入れかわるとき}\par\nobreak
光が、屈折率の大きい物質（光が遅く進む物質）との境で反射すると、反射した光は山と谷が入れかわります。経路が半波長ずれたのと同じです。屈折率の小さい物質との境での反射では、入れかわりません。これは、軽いひもから重いひもへ進んだ波が、境で山と谷を入れかえてはね返るのと同じしくみです（波-7 ⑤。重いひもほど波は遅い）。光でなぜそうなるのかは、物理シリーズの『電磁気学』（電-15）で、光の電場のつながり方から導きます。

\begin{bkcenter}
\begin{bktable}{colspec={Q[l,m,wd=42mm] Q[l,m,wd=48mm] Q[l,m,wd=34mm]}, rowsep=1.2pt, row{1}={font=\sffamily\small}, hlines={bkLine}, vline{2,3}={bkLine}}
反射する境 & 例 & 反射した光 \\
屈折率が小さい→大きい & 空気→水、空気→ガラス & 山と谷が入れかわる \\
屈折率が大きい→小さい & 水→空気、ガラス→空気 & そのまま \\
\end{bktable}
\bktabcaption{反射で山と谷が入れかわるか（光がどちらの物質から来て、境で反射するか）。}
\label{ph:波-13:tab-hansha}
\end{bkcenter}

\Needspace{5\baselineskip}\noindent\textbf{明るくなる条件、暗くなる条件}\par\nobreak
シャボン膜では、表の面で入れかわり（①）、裏の面ではそのまま（②）なので、2つの反射光は、光路長の差 $2nt$ のほかに、もう半波長ずれています。だから、ほぼ垂直に見るとき
\pheq{\text{強め合う：}\ 2nt=\Bigl(m+\frac12\Bigr)\lambda,\qquad \text{弱め合う：}\ 2nt=m\lambda\qquad(m=0,\ 1,\ 2,\ \cdots)}{m ＝ m}
です。入れかわりが2つの面で両方起きるとき（または両方起きないとき）は、ずれが打ち消し合うので、条件が逆になります（強め合う：$2nt=m\lambda$、弱め合う：$2nt=(m+\tfrac12)\lambda$）。条件を覚えるのではなく、「光路長の差」と「入れかわりの数」を1つずつ数えて決めましょう。

\Needspace{5\baselineskip}\noindent\textbf{ななめに見るとき：$2nt\cos r$}\par\nobreak
光が膜に角 $i$（入射角）で当たると、膜の中では屈折して角 $r$（屈折角）で進みます（$\sin i=n\sin r$、波-12 ③）。図\ref{ph:波-13:fig-naname}で、2つの光は、点 R から先は同じ道を通ります。それより前の、光路長の差を数えます。

\begin{bkfigure}
\begin{tikzpicture}[x=1mm, y=1mm]
  \fill[bkBlue!10] (0,8) rectangle (60,20);
  \draw[bkInk, line width=0.7pt] (0,20) -- (60,20); \draw[bkInk, line width=0.7pt] (0,8) -- (60,8);
  \coordinate (P) at (20,20); \coordinate (Q) at (27.12,8); \coordinate (R) at (34.25,20); \coordinate (H) at (25.9,27.0);
  \draw[rxn ray3] (9.28,29.0) -- (P);
  \draw[rxn ray3] (12.8,38.0) -- (R);
  \draw[bkBlue!70!black, line width=0.9pt] (P) -- (Q) -- (R);
  \draw[rxn ray] (R) -- (46.5,30.3);
  \draw[bkGray, line width=0.5pt, dashed] (P) -- (H);
  \draw[bkOrange!85!black, line width=1.6pt] (H) -- (R);
  \draw[bkGray, line width=0.4pt, dashed] (20,8) -- (20,33);
  \draw[bkGray, line width=0.4pt, dashed] (34.25,20) -- (34.25,35);
  \draw[bkInk, line width=0.5pt] (20,26) arc (90:140:6); \node[rxn lab] at (16.2,27.9) {$i$};
  \draw[bkInk, line width=0.5pt] (20,14) arc (-90:-59.3:6); \node[rxn lab] at (22.4,12.2) {$r$};
  \foreach \c/\n/\a in {P/P/north east, Q/Q/north, R/R/north west, H/H/south west}{\fill[bkInk] (\c) circle (0.6); \node[rxn lab, anchor=\a] at (\c) {\n};}
  \draw[rxn dim] (-2,8) -- (-2,20); \node[rxn lab, anchor=east] at (-2.5,14) {$t$};
  \node[rxn lab, anchor=west, align=left] at (64,32) {外の光の余分な道：太い線 HR\\HR $=$ PR $\sin i$};
  \node[rxn lab, anchor=west, align=left] at (64,16) {膜の中の道：PQ＋QR\\PQ $=$ QR $=t/\cos r$};
\end{tikzpicture}
\bkfigcaption{ななめに当たる光。左の光は P で膜に入り、Q で反射して R に出る。右の光は R で反射する。PH は2つの光に直角な線（同じ波面）で、そこから先の道のりを比べる。}
\label{ph:波-13:fig-naname}
\end{bkfigure}

膜の中の道 PQ＋QR は $2t/\cos r$ で、光路長は $2nt/\cos r$ です。外を通る光の余分な道 HR は、PR＝$2t\tan r$ を使って $\mathrm{HR}=2t\tan r\sin i$ です。$\sin i=n\sin r$ を入れて、差をとると
\[ \frac{2nt}{\cos r}-2t\tan r\cdot n\sin r=\frac{2nt}{\cos r}\bigl(1-\sin^2r\bigr)=2nt\cos r \]
となります（$\tan r=\sin r/\cos r$、$1-\sin^2r=\cos^2r$）。シャボン膜なら、反射の入れかわり1つを合わせて\label{ph:波-13:hakumaku}%
\pheq{\text{強め合う：}\ 2nt\cos r=\Bigl(m+\frac12\Bigr)\lambda}{m ＝ m}
です。ななめから見るほど $\cos r$ が小さくなり、強め合う波長が短い側へずれます。シャボン玉やチョウの羽の色が、見る角度で変わるのはこのためです。$i=0$ なら $r=0$、$\cos r=1$ で、垂直のときの式にもどります。

\Needspace{5\baselineskip}\noindent\textbf{シャボン玉の色と、黒い膜}\par\nobreak
白い光が当たると、膜の厚さ $t$ で決まる波長だけが強く反射し、その色に見えます（\nobreak\bkreidairef{5}\nobreak）。シャボン膜をたてにつるすと、液が下へ流れて、上ほど薄くなります。そのため、色の帯が横じまに並び、時間とともに下へ動いていきます。いちばん上が \qty{10}{nm} ほどまで薄くなると、$2nt$ は波長よりずっと小さく、ほぼ0になります。すると、反射の入れかわりのぶんだけ、どの色も打ち消し合い、膜は黒く見えます。黒くなった所は、もうすぐ割れる合図です。

\begin{phbaisuru}{膜の厚さ $t$}
$2nt$ が2倍になり、強め合う波長が2倍の数だけ現れます。グラフは、シャボン膜（$n=1.33$）の反射の明るさ（2つの光の重なり方。山を1）です。実線の \qty{300}{nm} では、強く反射するのは緑（\qty{532}{nm}）の1つの山だけですが、破線の \qty{600}{nm} では、青（\qty{456}{nm}）と赤（\qty{638}{nm}）の2つの山になります。もっと厚くすると、山がたくさん並んで、いろいろな色がまざり、白っぽく見えます。厚いガラスの板が色づかないのはこのためです。
\par\vspace{2pt}\centering
\begin{tikzpicture}
\begin{axis}[bkplot, width=104mm, height=42mm, xmin=380, xmax=780, ymin=0, ymax=1.25, xtick={400,450,500,550,600,650,700,750}, ytick={0,0.5,1}, xlabel={$\lambda$（nm）}, ylabel={反射の明るさ}]
  \addplot[bkcurve, domain=380:780, samples=200] {sin(180*2*1.33*300/x)^2};
  \addplot[bkcurve2, dashed, domain=380:780, samples=200] {sin(180*2*1.33*600/x)^2};
\end{axis}
\end{tikzpicture}
\end{phbaisuru}

\Needspace{5\baselineskip}\noindent\textbf{反射防止膜}\par\nobreak
めがねやカメラのレンズの表面では、ガラスに当たった光の約4\,\%が反射してしまいます（波-7 ⑨）。レンズを何枚も重ねたカメラでは、反射した光がレンズの中で行き来して、写真が白っぽくなります。そこで、ガラスの表面に、屈折率がガラスと空気の中間の薄い膜（たとえばフッ化マグネシウム、$n=1.38$）をつけます。これが\textbf{反射防止膜}です。\bkidx{はんしゃぼうしまく}{反射防止膜}

反射防止膜では、膜の表（空気→膜）でも、膜の裏（膜→ガラス、$1.38<1.52$）でも、入れかわりが起きます。2つの入れかわりは打ち消し合うので、弱め合う条件は $2nt=(m+\tfrac12)\lambda$ です。いちばん薄い $m=0$ では、膜の厚さを
\[ t=\frac{\lambda}{4n} \]
にします。緑の光（\qty{550}{nm}）に合わせると、$t=550\div(4\times1.38)\approx\qty{100}{nm}$ です（\nobreak\bkreidairef{6}\nobreak）。緑の反射が弱まり、残った赤と青の反射がまざるので、反射防止膜をつけたレンズは、反射が赤紫に見えます（何層も重ねた膜では、緑っぽく見えるものもあります）。

\begin{motto}[4]
\phboxhead{反射防止膜で、反射はどれだけ減る？}
2つの反射光の強さが違うと、打ち消しても0にはなりません。膜の両側の反射をすべて数える式（多重反射の式。専門の本で導く）で計算すると、ガラス（1.52）だけでは4.3\,\%の反射が、フッ化マグネシウムの膜（\qty{100}{nm}）をつけると、\qty{550}{nm} で1.3\,\%、\qty{400}{nm} で2.2\,\%、\qty{700}{nm} で1.6\,\%になります。反射を完全に0にするには、膜の屈折率を $\sqrt{1.52}\approx1.23$ にすればよいのですが、そのような丈夫な材料は少ないので、高級なレンズでは、屈折率の違う膜を何層も重ねて、広い波長の反射を0.5\,\%以下にしています。
\end{motto}

\Needspace{5\baselineskip}\noindent\textbf{くさび形の空気の層}\par\nobreak
平らなガラス板を2枚重ね、片方のはしに髪の毛を1本はさむと、2枚のあいだに、くさびの形の薄い空気の層ができます（図\ref{ph:波-13:fig-kusabi}）。上から単色の光を当てて見ると、髪の毛と平行な暗い線が、同じ間かくで並びます。

\begin{bkfigure}
\begin{tikzpicture}[x=1mm, y=1mm]
  \fill[rxn glass] (0,0) rectangle (100,4);
  \fill[rxn glass] (0,4) -- (100,10) -- (100,14) -- (0,8) -- cycle;
  \fill[bkInk!80] (97,7) circle (3);
  \node[rxn lab, anchor=west] at (101.5,4.5) {髪の毛（太さ $H$）};
  \foreach \x in {15,35,55,75}{\draw[rxn ray3] (\x,26) -- (\x,{8+0.06*\x+1});}
  \node[rxn lab, anchor=south] at (45,26.5) {単色の光（上から）};
  \draw[rxn dim] (0,-3) -- (97,-3); \node[rxn lab, fill=white] at (48.5,-3) {$L$};
  \draw[bkInk, line width=0.4pt] (62,4) -- (62,{4+0.06*62}); \draw[bkInk, line width=0.4pt] (61,4) -- (63,4); \draw[bkInk, line width=0.4pt] (61,{4+0.06*62}) -- (63,{4+0.06*62});
  \node[rxn lab, anchor=west] at (63.2,9.4) {$t$};
  \node[rxn lab, anchor=north west] at (0,-7) {上から見た暗い線};
  \begin{scope}[shift={(0,-22)}]
    \fill[bkYellow!35] (0,0) rectangle (100,10);
    \foreach \k in {0,...,9}{\fill[bkInk!75] ({10*\k-0.7},0) rectangle ({10*\k+0.7},10);}
    \draw[bkInk, line width=0.5pt] (0,0) rectangle (100,10);
    \draw[rxn dim] (30,11.5) -- (40,11.5); \node[rxn lab, anchor=south] at (35,12) {$\Delta x$};
  \end{scope}
\end{tikzpicture}
\bkfigcaption{くさび形の空気の層（上：横から。$t$ は空気の層の厚さで、大きく引きのばしてある。下：上から見た暗い線）。上のガラスの下の面（そのまま）と、下のガラスの上の面（入れかわる）で反射した光が干渉する。2枚がふれる左はしは暗い。}
\label{ph:波-13:fig-kusabi}
\end{bkfigure}

上のガラスの下の面での反射（ガラス→空気）はそのまま、下のガラスの上の面での反射（空気→ガラス）は入れかわるので、入れかわりは1つです。空気の屈折率は1なので、暗い線ができるのは
\[ 2t=m\lambda\qquad(m=0,\ 1,\ 2,\ \cdots) \]
の所です。ふれ合う左はし（$t=0$）は暗くなります。左はしから $x$ の所の厚さは、比例で $t=Hx/L$ です。暗い線の位置 $x_m$ は $2Hx_m/L=m\lambda$ から $x_m=mL\lambda/(2H)$ なので、となりどうしの間かくは\label{ph:波-13:kusabi}%
\pheq{\Delta x=\frac{L\lambda}{2H}}{m ＝ m × m ÷ m}
です。暗い線が1本進むごとに、空気の層の厚さが半波長（$\lambda/2$）ずつ増えている、と読めます。間かく $\Delta x$ をはかれば、髪の毛の太さ $H$ が、光の波長をものさしにしてはかれます（\nobreak\bkreidairef{7}\nobreak）。

\Needspace{5\baselineskip}\noindent\textbf{ニュートンリング}\par\nobreak
平らなガラス板の上に、片面がゆるやかにふくらんだレンズ（平凸レンズ）をのせると、レンズと板のあいだの空気の層は、中心から外へいくほど厚くなります。上から単色の光を当てると、同じ厚さの所が円になるので、明るい輪と暗い輪が同心円に並びます。これを\textbf{ニュートンリング}といいます。\bkidx{にゅーとんりんぐ}{ニュートンリング}ニュートンが詳しく調べたので、この名前があります。

\begin{bkfigure}
\begin{tikzpicture}[x=1mm, y=1mm]
  % 横から
  \fill[rxn glass] (-31,-4) rectangle (31,0);
  \fill[rxn glass] plot[domain=-28:28, samples=50] (\x,{80-sqrt(6400-\x*\x)}) -- (28,10) -- (-28,10) -- cycle;
  \draw[bkInk, line width=0.6pt] plot[domain=-28:28, samples=50] (\x,{80-sqrt(6400-\x*\x)});
  \draw[bkInk, line width=0.4pt] (20,0) -- (20,2.54); \draw[bkInk, line width=0.4pt] (19,2.54) -- (21,2.54);
  \node[rxn lab, anchor=west] at (21,1.3) {$t$};
  \draw[rxn dim] (0,-6.5) -- (20,-6.5); \node[rxn lab, anchor=north] at (10,-7) {$\rho$};
  \draw[bkGray, line width=0.4pt, dashed] (0,-4) -- (0,14);
  \node[rxn lab, anchor=south] at (0,14.3) {平凸レンズ};
  \node[rxn lab, anchor=north] at (-19,-5) {平らなガラス板};
  % 円の作図
  \begin{scope}[shift={(46,-6)}]
    \draw[bkGray, line width=0.5pt] (-15,4.02) arc (240:320:30);
    \coordinate (C) at (0,30); \coordinate (A) at (18,6); \coordinate (B) at (0,6);
    \draw[bkInk, line width=0.8pt] (C) -- (A) -- (B) -- cycle;
    \draw[bkInk, line width=0.8pt] (B) -- (0,0);
    \fill[bkInk] (C) circle (0.7); \node[rxn lab, anchor=south] at (0,31) {中心};
    \node[rxn lab, anchor=south west] at (10.5,19) {$R$};
    \node[rxn lab, anchor=east] at (-0.8,18) {$R-t$};
    \node[rxn lab, anchor=north] at (9,5.6) {$\rho$};
    \node[rxn lab, anchor=east] at (-0.8,3) {$t$};
  \end{scope}
  % 上から
  \begin{scope}[shift={(88,10)}]
    \fill[bkYellow!35] (0,0) circle (13);
    \foreach \m in {1,...,13}{\draw[bkInk!75, line width=0.7pt] (0,0) circle ({3.5*sqrt(\m)});}
    \fill[bkInk!80] (0,0) circle (1.4);
    \node[rxn lab, anchor=north] at (0,-14) {上から見た輪};
  \end{scope}
\end{tikzpicture}
\bkfigcaption{ニュートンリング。左：横から見た空気の層（厚さは大きく引きのばしてある）。中：レンズの球面の中心と、半径 $R$・$\rho$・$R-t$ の直角三角形。右：上から見た暗い輪。輪の半径は $\sqrt{m}$ に比例するので、外へいくほど間かくがせまくなる。}
\label{ph:波-13:fig-newton}
\end{bkfigure}

レンズの球面の半径を $R$、中心から $\rho$（ロー）の所の空気の層の厚さを $t$ とします（輪の半径は、屈折角 $r$ とぶつけないように $\rho$ と書きます）。図\ref{ph:波-13:fig-newton}の直角三角形で、三平方の定理（4-4）から
\[ (R-t)^2+\rho^2=R^2\qquad\text{より}\qquad \rho^2=2Rt-t^2 \]
です。$t$ は数 \unit{\micro m} で、$R$（約 \qty{1}{m}）よりずっと小さいので、$t^2$ を無視して $t\approx \rho^2/(2R)$ です。反射の入れかわりは、くさびと同じく1つなので、暗い輪は $2t=m\lambda$ の所にできます。$2\times \rho^2/(2R)=m\lambda$ から
\pheq{\rho_m^{\,2}=m\lambda R\qquad(m=0,\ 1,\ 2,\ \cdots)}{m$^2$ ＝ m × m}
です。中心（$m=0$）は暗く、$m$ 番めの暗い輪の半径は $\sqrt{m}$ に比例します。輪の半径をはかると、レンズの球面の半径 $R$ がわかります（\nobreak\bkreidairef{8}\nobreak）。

\begin{phmijika}[光で平らさを調べる]
レンズや鏡を作る工場では、とても平らに仕上げた基準のガラス（オプチカルフラット）を、調べる面にのせ、単色の光で干渉じまを見ます。しまがまっすぐ等間かくなら面は平ら、曲がっていれば、そこにでこぼこがあります。しま1本が厚さの $\lambda/2$（約 \qty{0.3}{\micro m}）にあたるので、ふれずに、光の波長の数分の1の精度で形が調べられます。
\end{phmijika}

\begin{phmada}[シャボン膜は、いつ割れる？]
シャボン膜が黒くなってから割れることは、④の干渉で説明できます。けれども、ある膜が「何秒後に割れるか」を前もって言い当てることは、今もできません。膜の液が重さで下へ流れて薄くなる速さは、液の粘り気や、表面の石けんの分子の動きで決まり、ある程度は計算できます。ところが実際には、空気の流れ、水の蒸発、目に見えないちりなどで、膜の一部が急に薄くなって割れ、同じ条件で作っても、割れるまでの時間が大きくばらつきます。膜が薄くなっていくようすを干渉の色で細かくはかりながら、割れるしくみを調べる研究が、今も続いています。
\end{phmada}

\phkokomade

\Needspace{12\baselineskip}
\phstage{shiki}{波の式とゆれのモデルで調べる}{5}

\bksubsection{5}{円い穴の回折と、見分けられる限界}

\leavevmode\label{ph:波-13:sec5}目・カメラ・望遠鏡・顕微鏡は、どれも円い穴（ひとみ・レンズの口）から光を取りこみます。光は波なので、円い穴を通るときにも回折して広がります。そのため、星のような1つの点の像も、小さな点にはならず、少し広がった円になります。この広がりが、細かい物を見分けられる限界（\textbf{分解能}）を決めます。\bkidx{ぶんかいのう}{分解能}

\Needspace{5\baselineskip}\noindent\textbf{すきまの回折を、フーリエ変換で見る}\par\nobreak
まず、はば $a$ の1つのすきまを思い出します（波-8 ④）。すきまの中の各点から出た波を、遠くで（角 $\theta$ の向きで）たし合わせると、はじめて暗くなる向きは $\sin\theta=\lambda/a$ でした。

これを、すきまの中を細かく分けた「矢印のたし算」として見直します。すきまの中の高さ $y$ の点から出た光は、まん中の点の光より、経路が $y\sin\theta$ だけ短く、位相が $2\pi y\sin\theta/\lambda$ だけ進んでいます。$q=2\pi\sin\theta/\lambda$ とおくと（9-4 の $\omega$ にあたる量です。光の角振動数 $\omega$ とぶつけないように $q$ と書きます）、すべての点の矢印の和は（進むかおくれるかの向きは、和の大きさに効かないので、9-4 と同じ形に書きます）
\[ \int_{-a/2}^{a/2}e^{-jqy}\,\dd y \]
です。これは、はば $a$ の方形パルスのフーリエ変換（9-4④）と同じ形です。9-4④の計算で、はばを2から $a$ に、$\omega$ を $q$ に変えると、答えは $2\sin(qa/2)/q$ になり、はじめて0になるのは $qa/2=\pi$、つまり $\sin\theta=\lambda/a$ です。波-8 ④と同じ答えが出ました。

\begin{phyomi}[遠くで見る回折の模様＝穴の形のフーリエ変換]
「遠くで見る光の広がり方は、穴の形を変換した形になる」。9-4④の「時間で短いほど、周波数で広い」は、ここでは「すきまがせまいほど、光は広い角に広がる」と読めます。穴の形さえわかれば、回折の模様は計算で決まります。
\end{phyomi}

\Needspace{5\baselineskip}\noindent\textbf{円い穴では 1.22 倍}\par\nobreak
直径 $D$ の円い穴でも、同じように穴の中の矢印をすべてたします。円は、まん中の高さでははばが $D$ ですが、上と下にいくほどせまくなります。はばの平均が $D$ より小さいぶん、すきまより少し広く広がります。計算すると、点の像は、明るい円板とそれを囲む薄い輪になり（図\ref{ph:波-13:fig-airy}）、はじめて暗くなる向きは
\pheq{\sin\theta=1.22\,\frac{\lambda}{D}}{（単位なし）＝ m ÷ m}
\label{ph:波-13:enkou}です。中心の明るい円板を、エアリーの円板といいます。1.22 という数は、\textbf{ベッセル関数}という関数が0になる所から決まります。ベッセル関数の作り方と性質は専門の本で学びます。ここでは、計算した値を表で使います。

\begin{bkcenter}
\begin{bktable}{colspec={Q[l,m,wd=44mm] Q[c,m,wd=36mm] Q[c,m,wd=36mm]}, rowsep=1.2pt, row{1}={font=\sffamily\small}, hlines={bkLine}, vline{2,3}={bkLine}}
 & すきま（はば $a$） & 円い穴（直径 $D$） \\
1つめの暗い所 & 1.00 & 1.22 \\
1つめの明るい所（強さ） & 1.43（4.7\,\%） & 1.63（1.7\,\%） \\
2つめの暗い所 & 2.00 & 2.23 \\
2つめの明るい所（強さ） & 2.46（1.6\,\%） & 2.68（0.4\,\%） \\
3つめの暗い所 & 3.00 & 3.24 \\
\end{bktable}
\bktabcaption{回折の暗い所・明るい所の向き。数は $\sin\theta$ を、すきまでは $\lambda/a$、円い穴では $\lambda/D$ を単位にして表したもの。強さは、まん中を100\,\%とした値。円い穴の数は、ベッセル関数 $J_1$ から計算した（たとえば 1.22 は、$J_1$ がはじめて0になる 3.832 を $\pi$ でわった値）。}
\label{ph:波-13:tab-bessel}
\end{bkcenter}

\begin{bkfigure}
\begin{tikzpicture}
\begin{axis}[bkplot, width=84mm, height=44mm, xmin=0, xmax=3.5, ymin=0, ymax=1.15, xtick={0,1,2,3}, ytick={0,0.5,1}, xlabel={$\sin\theta$}, ylabel={明るさ（まん中を1）}]
  \addplot[bkcurve2, dashed] coordinates {(0.000,1.0000) (0.050,0.9918) (0.100,0.9675) (0.150,0.9281) (0.200,0.8751) (0.250,0.8106) (0.300,0.7368) (0.350,0.6566) (0.400,0.5728) (0.450,0.4881) (0.500,0.4053) (0.550,0.3267) (0.600,0.2546) (0.650,0.1904) (0.700,0.1353) (0.750,0.0901) (0.800,0.0547) (0.850,0.0289) (0.900,0.0119) (0.950,0.0027) (1.000,0.0000) (1.050,0.0022) (1.100,0.0080) (1.150,0.0158) (1.200,0.0243) (1.250,0.0324) (1.300,0.0392) (1.350,0.0441) (1.400,0.0468) (1.450,0.0470) (1.500,0.0450) (1.550,0.0411) (1.600,0.0358) (1.650,0.0295) (1.700,0.0229) (1.750,0.0165) (1.800,0.0108) (1.850,0.0061) (1.900,0.0027) (1.950,0.0007) (2.000,0.0000) (2.050,0.0006) (2.100,0.0022) (2.150,0.0045) (2.200,0.0072) (2.250,0.0100) (2.300,0.0125) (2.350,0.0146) (2.400,0.0159) (2.450,0.0165) (2.500,0.0162) (2.550,0.0152) (2.600,0.0136) (2.650,0.0115) (2.700,0.0091) (2.750,0.0067) (2.800,0.0045) (2.850,0.0026) (2.900,0.0012) (2.950,0.0003) (3.000,0.0000) (3.050,0.0003) (3.100,0.0010) (3.150,0.0021) (3.200,0.0034) (3.250,0.0048) (3.300,0.0061) (3.350,0.0072) (3.400,0.0079) (3.450,0.0083) (3.500,0.0083)};
  \addplot[bkcurve] coordinates {(0.000,1.0000) (0.050,0.9938) (0.100,0.9756) (0.150,0.9458) (0.200,0.9053) (0.250,0.8553) (0.300,0.7975) (0.350,0.7332) (0.400,0.6645) (0.450,0.5931) (0.500,0.5209) (0.550,0.4495) (0.600,0.3806) (0.650,0.3157) (0.700,0.2558) (0.750,0.2018) (0.800,0.1544) (0.850,0.1139) (0.900,0.0803) (0.950,0.0534) (1.000,0.0328) (1.050,0.0180) (1.100,0.0081) (1.150,0.0025) (1.200,0.0002) (1.250,0.0004) (1.300,0.0023) (1.350,0.0052) (1.400,0.0085) (1.450,0.0116) (1.500,0.0143) (1.550,0.0162) (1.600,0.0173) (1.650,0.0175) (1.700,0.0168) (1.750,0.0154) (1.800,0.0135) (1.850,0.0113) (1.900,0.0090) (1.950,0.0067) (2.000,0.0046) (2.050,0.0028) (2.100,0.0015) (2.150,0.0005) (2.200,0.0001) (2.250,0.0000) (2.300,0.0003) (2.350,0.0008) (2.400,0.0015) (2.450,0.0022) (2.500,0.0029) (2.550,0.0035) (2.600,0.0039) (2.650,0.0041) (2.700,0.0041) (2.750,0.0040) (2.800,0.0036) (2.850,0.0031) (2.900,0.0026) (2.950,0.0020) (3.000,0.0014) (3.050,0.0009) (3.100,0.0005) (3.150,0.0002) (3.200,0.0000) (3.250,0.0000) (3.300,0.0001) (3.350,0.0003) (3.400,0.0005) (3.450,0.0008) (3.500,0.0010)};
  \node[font=\sffamily\footnotesize, anchor=west, align=left, fill=white, inner sep=1.5pt] at (axis cs:1.55,0.8) {実線：円い穴\\破線：すきま};
  \draw[bkInk, line width=0.5pt, -{Stealth[length=2.4pt]}] (axis cs:1.22,0.32) -- (axis cs:1.22,0.04);
  \node[font=\sffamily\footnotesize, anchor=south] at (axis cs:1.22,0.32) {1.22};
\end{axis}
\end{tikzpicture}\hspace{3mm}%
\begin{tikzpicture}[x=0.85mm, y=0.85mm]
  \fill[bkInk] (-13,-13) rectangle (13,13);
  \shade[inner color=white, outer color=bkInk] (0,0) circle (5.2);
  \draw[white!55!bkInk, line width=0.9pt] (0,0) circle (7.0);
  \draw[white!30!bkInk, line width=0.7pt] (0,0) circle (11.4);
  \node[rxn lab, anchor=north] at (0,-14) {円い穴の像};
\end{tikzpicture}
\bkfigcaption{左：すきま（破線）と円い穴（実線）の回折の明るさ（横軸の $\sin\theta$ は、すきまでは $\lambda/a$、円い穴では $\lambda/D$ が単位）。円い穴のほうが、はじめて暗くなる所が 1.22 倍外にある。右：円い穴を通った光の、点の像（エアリーの円板と、薄い輪）。輪は強さが数\,\%しかないので、図では見やすくかいた。}
\label{ph:波-13:fig-airy}
\end{bkfigure}

\Needspace{5\baselineskip}\noindent\textbf{2つの点を見分ける：レイリーの基準}\par\nobreak
夜空の2つの星が近くに並んでいるとします。それぞれの星の像は、エアリーの円板に広がります。2つの星の向きの角が小さいと、2つの円板が重なって、1つの星に見えてしまいます。そこで、「一方の像の中心が、もう一方の像のはじめの暗い輪の上にくるとき、ちょうど見分けられる」と決めます。これを\textbf{レイリーの基準}といいます（1879年）。このとき、2つの像を合わせた明るさは、まん中で約74\,\%にへこみ、2つの山として見えます（図\ref{ph:波-13:fig-rayleigh}）。角は小さいので $\sin\theta\approx\theta$（5-17⑥）として、見分けられる最小の角は
\pheq{\theta_{\min}\approx1.22\,\frac{\lambda}{D}\quad(\mathrm{rad})}{rad ＝ m ÷ m}
\label{ph:波-13:rayleigh-kijun}です。

\begin{bkfigure}
\begin{tikzpicture}
\begin{axis}[bkplot, axis lines=left, width=110mm, height=46mm, xmin=-1.5, xmax=2.72, ymin=0, ymax=1.18, xtick={-1,0,1,2}, xticklabels={$-1$,0,1,2}, ytick={0,0.5,1}, xlabel={$\theta$（$\lambda/D$ が単位）}, ylabel={明るさ}]
  \addplot[bkthin, dashed] coordinates {(-1.500,0.0143) (-1.450,0.0116) (-1.400,0.0084) (-1.349,0.0051) (-1.299,0.0022) (-1.249,0.0003) (-1.199,0.0002) (-1.148,0.0026) (-1.098,0.0084) (-1.048,0.0185) (-0.998,0.0337) (-0.947,0.0546) (-0.897,0.0820) (-0.847,0.1162) (-0.797,0.1574) (-0.746,0.2055) (-0.696,0.2601) (-0.646,0.3208) (-0.596,0.3864) (-0.545,0.4559) (-0.495,0.5277) (-0.445,0.6003) (-0.395,0.6719) (-0.345,0.7405) (-0.294,0.8044) (-0.244,0.8617) (-0.194,0.9108) (-0.144,0.9502) (-0.093,0.9787) (-0.043,0.9954) (0.007,0.9999) (0.057,0.9919) (0.108,0.9718) (0.158,0.9401) (0.208,0.8978) (0.258,0.8462) (0.309,0.7868) (0.359,0.7214) (0.409,0.6517) (0.459,0.5797) (0.510,0.5071) (0.560,0.4358) (0.610,0.3673) (0.660,0.3030) (0.710,0.2439) (0.761,0.1911) (0.811,0.1449) (0.861,0.1058) (0.911,0.0735) (0.962,0.0480) (1.012,0.0288) (1.062,0.0152) (1.112,0.0064) (1.163,0.0016) (1.213,0.0000) (1.263,0.0007) (1.313,0.0030) (1.364,0.0061) (1.414,0.0094) (1.464,0.0124) (1.514,0.0149) (1.565,0.0166) (1.615,0.0174) (1.665,0.0173) (1.715,0.0165) (1.765,0.0149) (1.816,0.0129) (1.866,0.0106) (1.916,0.0082) (1.966,0.0059) (2.017,0.0039) (2.067,0.0023) (2.117,0.0011) (2.167,0.0003) (2.218,0.0000) (2.268,0.0001) (2.318,0.0005) (2.368,0.0010) (2.419,0.0017) (2.469,0.0025) (2.519,0.0031) (2.569,0.0037) (2.620,0.0040) (2.670,0.0042) (2.720,0.0041)};
  \addplot[bkthin, dashed] coordinates {(-1.500,0.0041) (-1.450,0.0042) (-1.400,0.0040) (-1.349,0.0037) (-1.299,0.0031) (-1.249,0.0025) (-1.199,0.0017) (-1.148,0.0010) (-1.098,0.0005) (-1.048,0.0001) (-0.998,0.0000) (-0.947,0.0003) (-0.897,0.0011) (-0.847,0.0023) (-0.797,0.0039) (-0.746,0.0059) (-0.696,0.0082) (-0.646,0.0106) (-0.596,0.0129) (-0.545,0.0149) (-0.495,0.0165) (-0.445,0.0173) (-0.395,0.0174) (-0.345,0.0166) (-0.294,0.0149) (-0.244,0.0124) (-0.194,0.0094) (-0.144,0.0061) (-0.093,0.0030) (-0.043,0.0007) (0.007,0.0000) (0.057,0.0016) (0.108,0.0064) (0.158,0.0152) (0.208,0.0288) (0.258,0.0480) (0.309,0.0735) (0.359,0.1058) (0.409,0.1449) (0.459,0.1911) (0.510,0.2439) (0.560,0.3030) (0.610,0.3673) (0.660,0.4358) (0.710,0.5071) (0.761,0.5797) (0.811,0.6517) (0.861,0.7214) (0.911,0.7868) (0.962,0.8462) (1.012,0.8978) (1.062,0.9401) (1.112,0.9718) (1.163,0.9919) (1.213,0.9999) (1.263,0.9954) (1.313,0.9787) (1.364,0.9502) (1.414,0.9108) (1.464,0.8617) (1.514,0.8044) (1.565,0.7405) (1.615,0.6719) (1.665,0.6003) (1.715,0.5277) (1.765,0.4559) (1.816,0.3864) (1.866,0.3208) (1.916,0.2601) (1.966,0.2055) (2.017,0.1574) (2.067,0.1162) (2.117,0.0820) (2.167,0.0546) (2.218,0.0337) (2.268,0.0185) (2.318,0.0084) (2.368,0.0026) (2.419,0.0002) (2.469,0.0003) (2.519,0.0022) (2.569,0.0051) (2.620,0.0084) (2.670,0.0116) (2.720,0.0143)};
  \addplot[bkcurve] coordinates {(-1.500,0.0184) (-1.450,0.0158) (-1.400,0.0124) (-1.349,0.0088) (-1.299,0.0054) (-1.249,0.0028) (-1.199,0.0020) (-1.148,0.0037) (-1.098,0.0089) (-1.048,0.0186) (-0.998,0.0337) (-0.947,0.0550) (-0.897,0.0831) (-0.847,0.1185) (-0.797,0.1613) (-0.746,0.2114) (-0.696,0.2683) (-0.646,0.3313) (-0.596,0.3993) (-0.545,0.4708) (-0.495,0.5442) (-0.445,0.6177) (-0.395,0.6893) (-0.345,0.7572) (-0.294,0.8193) (-0.244,0.8742) (-0.194,0.9202) (-0.144,0.9563) (-0.093,0.9817) (-0.043,0.9962) (0.007,0.9999) (0.057,0.9935) (0.108,0.9781) (0.158,0.9552) (0.208,0.9266) (0.258,0.8943) (0.309,0.8604) (0.359,0.8272) (0.409,0.7967) (0.459,0.7708) (0.510,0.7511) (0.560,0.7388) (0.610,0.7346) (0.660,0.7388) (0.710,0.7511) (0.761,0.7708) (0.811,0.7967) (0.861,0.8272) (0.911,0.8604) (0.962,0.8943) (1.012,0.9266) (1.062,0.9552) (1.112,0.9781) (1.163,0.9935) (1.213,0.9999) (1.263,0.9962) (1.313,0.9817) (1.364,0.9563) (1.414,0.9202) (1.464,0.8742) (1.514,0.8193) (1.565,0.7572) (1.615,0.6893) (1.665,0.6177) (1.715,0.5442) (1.765,0.4708) (1.816,0.3993) (1.866,0.3313) (1.916,0.2683) (1.966,0.2114) (2.017,0.1613) (2.067,0.1185) (2.117,0.0831) (2.167,0.0550) (2.218,0.0337) (2.268,0.0186) (2.318,0.0089) (2.368,0.0037) (2.419,0.0020) (2.469,0.0028) (2.519,0.0054) (2.569,0.0088) (2.620,0.0124) (2.670,0.0158) (2.720,0.0184)};
  \draw[bkOrange!85!black, line width=0.6pt, {Stealth[length=2.6pt]}-{Stealth[length=2.6pt]}] (axis cs:0.61,0.74) -- (axis cs:0.61,1.0);
  \node[font=\sffamily\footnotesize, anchor=west] at (axis cs:1.85,0.6) {実線：2つの和};
  \node[font=\sffamily\footnotesize, anchor=west] at (axis cs:1.85,0.42) {破線：1つずつ};
  \node[font=\sffamily\footnotesize, anchor=west] at (axis cs:1.8,0.86) {へこみ（74\,\%）};
  \node[font=\sffamily\footnotesize, anchor=south] at (axis cs:1.22,1.02) {1.22};
\end{axis}
\end{tikzpicture}
\bkfigcaption{レイリーの基準。2つの点の像（破線）の間の角が $1.22\lambda/D$ のとき、合わせた明るさ（実線）は、まん中で山の約74\,\%にへこみ、2つの点として見分けられる。角を $0.8\lambda/D$ にすると、へこみがなくなり、1つの山になる。}
\label{ph:波-13:fig-rayleigh}
\end{bkfigure}

\begin{phbaisuru}{穴の直径 $D$}
見分けられる最小の角は半分になります。グラフの実線は直径 $D$、破線は直径 $2D$ の穴の、点の像の明るさです（横軸は、直径 $D$ の $\lambda/D$ を単位にした角）。$2D$ では、はじめて暗くなる角が 1.22 から 0.61 に縮みます。大きな望遠鏡ほど、遠くの星を細かく見分けられるのはこのためです。
\par\vspace{2pt}\centering
\begin{tikzpicture}
\begin{axis}[bkplot, width=104mm, height=40mm, xmin=0, xmax=2.5, ymin=0, ymax=1.15, xtick={0,0.61,1.22,2}, xticklabels={0,0.61,1.22,2}, ytick={0,0.5,1}, xlabel={角（$\lambda/D$ が単位）}, ylabel={明るさ}]
  \addplot[bkcurve] coordinates {(0.000,1.0000) (0.042,0.9957) (0.083,0.9830) (0.125,0.9621) (0.167,0.9334) (0.208,0.8976) (0.250,0.8553) (0.292,0.8076) (0.333,0.7552) (0.375,0.6993) (0.417,0.6409) (0.458,0.5811) (0.500,0.5209) (0.542,0.4613) (0.583,0.4032) (0.625,0.3476) (0.667,0.2951) (0.708,0.2463) (0.750,0.2018) (0.792,0.1618) (0.833,0.1266) (0.875,0.0962) (0.917,0.0706) (0.958,0.0495) (1.000,0.0328) (1.042,0.0201) (1.083,0.0109) (1.125,0.0048) (1.167,0.0014) (1.208,0.0001) (1.250,0.0004) (1.292,0.0019) (1.333,0.0041) (1.375,0.0068) (1.417,0.0095) (1.458,0.0121) (1.500,0.0143) (1.542,0.0160) (1.583,0.0170) (1.625,0.0175) (1.667,0.0173) (1.708,0.0166) (1.750,0.0154) (1.792,0.0139) (1.833,0.0121) (1.875,0.0102) (1.917,0.0082) (1.958,0.0063) (2.000,0.0046) (2.042,0.0031) (2.083,0.0019) (2.125,0.0009) (2.167,0.0003) (2.208,0.0000) (2.250,0.0000) (2.292,0.0002) (2.333,0.0006) (2.375,0.0011) (2.417,0.0017) (2.458,0.0023) (2.500,0.0029)};
  \addplot[bkcurve2, dashed] coordinates {(0.000,1.0000) (0.025,0.9938) (0.050,0.9756) (0.075,0.9458) (0.100,0.9053) (0.125,0.8553) (0.150,0.7975) (0.175,0.7332) (0.200,0.6645) (0.225,0.5931) (0.250,0.5209) (0.275,0.4495) (0.300,0.3806) (0.325,0.3157) (0.350,0.2558) (0.375,0.2018) (0.400,0.1544) (0.425,0.1139) (0.450,0.0803) (0.475,0.0534) (0.500,0.0328) (0.525,0.0180) (0.550,0.0081) (0.575,0.0025) (0.600,0.0002) (0.625,0.0004) (0.650,0.0023) (0.675,0.0052) (0.700,0.0085) (0.725,0.0116) (0.750,0.0143) (0.775,0.0162) (0.800,0.0173) (0.825,0.0175) (0.850,0.0168) (0.875,0.0154) (0.900,0.0135) (0.925,0.0113) (0.950,0.0090) (0.975,0.0067) (1.000,0.0046) (1.025,0.0028) (1.050,0.0015) (1.075,0.0005) (1.100,0.0001) (1.125,0.0000) (1.150,0.0003) (1.175,0.0008) (1.200,0.0015) (1.225,0.0022) (1.250,0.0029) (1.275,0.0035) (1.300,0.0039) (1.325,0.0041) (1.350,0.0041) (1.375,0.0040) (1.400,0.0036) (1.425,0.0031) (1.450,0.0026) (1.475,0.0020) (1.500,0.0014) (1.525,0.0009) (1.550,0.0005) (1.575,0.0002) (1.600,0.0000) (1.625,0.0000) (1.650,0.0001) (1.675,0.0003) (1.700,0.0005) (1.725,0.0008) (1.750,0.0010) (1.775,0.0013) (1.800,0.0014) (1.825,0.0016) (1.850,0.0016) (1.875,0.0016) (1.900,0.0014) (1.925,0.0013) (1.950,0.0011) (1.975,0.0008) (2.000,0.0006) (2.025,0.0004) (2.050,0.0002) (2.075,0.0001) (2.100,0.0000) (2.125,0.0000) (2.150,0.0000) (2.175,0.0001) (2.200,0.0002) (2.225,0.0003) (2.250,0.0005) (2.275,0.0006) (2.300,0.0007) (2.325,0.0008) (2.350,0.0008) (2.375,0.0008) (2.400,0.0007) (2.425,0.0006) (2.450,0.0005) (2.475,0.0004) (2.500,0.0003)};
\end{axis}
\end{tikzpicture}
\end{phbaisuru}

\begin{phmijika}[目・顕微鏡・ブルーレイ]
人のひとみは、明るい所で直径約 \qty{3}{mm} です。緑の光では $\theta_{\min}\approx2.2\times10^{-4}$\,rad（約0.8分。1分は1度の60分の1）で、視力1.0の人が見分ける約1分に近い値です（\nobreak\bkreidairef{9}\nobreak）。光の顕微鏡で見分けられる大きさも、同じ理由で、光の波長の半分ほど（約 \qty{0.2}{\micro m}）が限界です。それより小さいウイルスなどを見るには、波長のとても短い電子の波を使う電子顕微鏡を使います（原-3）。光ディスクでは、CD（\qty{780}{nm}）→ DVD（\qty{650}{nm}）→ ブルーレイ（\qty{405}{nm}）と、読み取る光の波長を短くし、レンズの口を大きくすることで、光の点を小さくしぼり、より細かい点の列を読めるようにしてきました。
\end{phmijika}

\Needspace{10\baselineskip}
\begin{phflow}[回折の広がり]
\phflowrow{すきまの各点の\\波の和（フーリエ変換）}{はば $a$ のすきま\\遠くで見る}{$\sin\theta=\dfrac{\lambda}{a}$（波-8 ④）}
\phflowrow{穴の各点の波の和}{直径 $D$ の円い穴\\小さな角}{$\theta_{\min}\approx1.22\,\dfrac{\lambda}{D}$}
\end{phflow}

\bksubsection{6}{偏光：光のゆれの向き}

\leavevmode\label{ph:波-13:sec6}波には、ゆれの向きが進む向きに垂直な横波と、平行な縦波があります（波-5 ①）。光は横波です。光の電場（詳しくは電-15）は、光の進む向きに垂直な面の中でゆれています。太陽や電球の光（\textbf{自然光}）は、ゆれの向きが1秒に何兆回も入れかわり、あらゆる向きのゆれがまざっています。ゆれの向きがそろった光を\textbf{偏光}といいます。\bkidx{へんこう}{偏光}

\Needspace{5\baselineskip}\noindent\textbf{偏光板}\par\nobreak
\label{ph:波-13:henkou}\textbf{偏光板}は、ある1つの向き（通す向き）のゆれだけを通し、それに垂直なゆれを吸いとる板です。ひもの波を、たてのすきまのある板に通すと、たてのゆれは通りぬけ、横のゆれは止められます。それと同じことが、光の電場のゆれで起きます。偏光板を2枚重ねて回すと、明るさが変わるのは、光が横波だからです。縦波（音）では、ゆれの向きは進む向きと決まっているので、このようなことは起きません。

\begin{bkfigure}
\begin{tikzpicture}[x=1mm, y=1mm]
  \draw[bkGray, line width=0.5pt, -{Stealth[length=3pt]}] (0,10) -- (122,10);
  % 自然光
  \begin{scope}[shift={(10,10)}]
    \draw[bkInk!50, line width=0.4pt] (0,0) ellipse (2.6 and 7.5);
    \foreach \a in {0,30,...,150}{\draw[bkRed!80!black, line width=0.8pt, {Stealth[length=2pt]}-{Stealth[length=2pt]}] ({-2.2*sin(\a)},{-6.5*cos(\a)}) -- ({2.2*sin(\a)},{6.5*cos(\a)});}
  \end{scope}
  \node[rxn lab, align=center, anchor=north] at (10,0.5) {自然光\\（いろいろな向き）};
  \rxnpol{(36,10)}{0}
  \node[rxn lab, align=center, anchor=north] at (36,0.5) {偏光板1\\（通す向き：たて）};
  \begin{scope}[shift={(60,10)}]
    \draw[bkInk!50, line width=0.4pt] (0,0) ellipse (2.6 and 7.5);
    \draw[bkRed!80!black, line width=1pt, {Stealth[length=2.4pt]}-{Stealth[length=2.4pt]}] (0,-6.5) -- (0,6.5);
  \end{scope}
  \node[rxn lab, align=center, anchor=north] at (60,0.5) {たてにゆれる光\\（強さ $\tfrac12$）};
  \rxnpol{(86,10)}{-60}
  \node[rxn lab, align=center, anchor=north] at (86,0.5) {偏光板2\\（$\theta$ だけ回す）};
  \begin{scope}[shift={(110,10)}]
    \draw[bkInk!50, line width=0.4pt] (0,0) ellipse (2.6 and 7.5);
    \draw[bkRed!80!black, line width=1pt, {Stealth[length=2.4pt]}-{Stealth[length=2.4pt]}] (-1.4,-2.4) -- (1.4,2.4);
  \end{scope}
  \node[rxn lab, align=center, anchor=north] at (110,0.5) {強さ $\tfrac12\cos^2\theta$};
\end{tikzpicture}
\bkfigcaption{偏光板を2枚通る光。だ円の中の両向きの矢印が、進む向きに垂直な面の中の光のゆれ。1枚めでたてのゆれだけが残り、2枚めは、通す向きの成分だけを通す。}
\label{ph:波-13:fig-henkou}
\end{bkfigure}

\Needspace{5\baselineskip}\noindent\textbf{マリュスの法則}\par\nobreak
\noindent\begin{minipage}[c]{0.66\linewidth}\raggedright
たてにゆれる光（振幅 $E_0$）が、通す向きが角 $\theta$ だけ傾いた偏光板に来たとします。ゆれを、通す向きの成分 $E_0\cos\theta$ と、それに垂直な成分 $E_0\sin\theta$ に分けます（ベクトルの分解、4-11）。偏光板は前者だけを通します。光の強さは振幅の2乗に比例するので、通りぬけた強さは
\end{minipage}\hfill
\begin{minipage}[c]{0.3\linewidth}\centering
\begin{tikzpicture}[x=1mm, y=1mm]
  \draw[bkGray, line width=0.5pt, dashed] (0,-2) -- (0,24);
  \draw[bkGray, line width=0.5pt, dashed] (-2.5,-1.4) -- ({24*sin(30)},{24*cos(30)});
  \draw[bkRed!80!black, line width=1.2pt, -{Stealth[length=3pt]}] (0,0) -- (0,20);
  \draw[bkBlue!80!black, line width=1.2pt, -{Stealth[length=3pt]}] (0,0) -- ({17.32*sin(30)},{17.32*cos(30)});
  \draw[bkGray, line width=0.5pt, dotted] (0,20) -- ({17.32*sin(30)},{17.32*cos(30)});
  \draw[bkInk, line width=0.5pt] (0,7) arc (90:60:7); \node[rxn lab] at (1.9,9.1) {$\theta$};
  \node[rxn lab, anchor=east] at (-0.8,15) {$E_0$};
  \node[rxn lab, anchor=west, align=left] at (9.5,13) {$E_0\cos\theta$\\（通る）};
  \node[rxn lab, anchor=south] at (12,24.2) {通す向き};
\end{tikzpicture}
\end{minipage}
\pheq{I=I_0\cos^2\theta}{W/m$^2$ ＝ W/m$^2$}
\label{ph:波-13:malus}となります。$I_0$ は、2枚めの偏光板に来る光の強さです。これを\textbf{マリュスの法則}といいます（マリュス、1809年ごろ）。\bkidx{まりゅすのほうそく}{マリュスの法則}自然光を1枚めの偏光板に通すと、あらゆる向きの $\cos^2$ を平均することになり、強さは $\tfrac12$ になります（$\cos^2$ の平均は $\tfrac12$：4-9⑧）。

\begin{phyomi}[$I=I_0\cos^2\theta$]
「通す向きとゆれの向きが同じ（$\theta=0$）なら全部、直角（$\theta=90°$）なら0。そのあいだは $\cos$ の2乗の割合」。2枚の偏光板を直角に重ねると、光はまったく通りません。
\end{phyomi}
\begin{phbaisuru}{角 $\theta$（$30°$ を $60°$ に）}
$\cos^230°=0.75$ から $\cos^260°=0.25$ へ、通る光は3分の1になります。角に比例して減るのではなく、$0°$ の近くではゆっくり、$45°$ の近くで速く減り、$90°$ でちょうど0になります。
\par\vspace{2pt}\centering
\begin{tikzpicture}
\begin{axis}[bkplot, width=104mm, height=40mm, xmin=0, xmax=180, ymin=0, ymax=1.15, xtick={0,30,60,90,120,150,180}, ytick={0,0.25,0.5,0.75,1}, xlabel={$\theta$（°）}, ylabel={$I/I_0$}]
  \addplot[bkcurve, domain=0:180, samples=91] {cos(x)^2};
  \addplot[only marks, mark=*, mark size=1.8pt, bkOrange] coordinates {(30,0.75) (60,0.25)};
\end{axis}
\end{tikzpicture}
\end{phbaisuru}

\begin{phmijika}[液晶の画面]
スマホやパソコンの液晶の画面では、光源の光をまず偏光板でそろえ、液晶の層で偏光の向きを回してから、通す向きが直角のもう1枚の偏光板に通します。液晶にかける電圧を変えると、偏光の向きの回り方が変わり、2枚めを通る光の量が変わります。こうして、点ごとに明るさを変えています（液晶の種類によって、しくみには細かい違いがあります）。やってみようの④で、偏光サングラスを回すと画面が暗くなったのは、画面の光がもともと偏光しているからです。
\end{phmijika}

\Needspace{5\baselineskip}\noindent\textbf{反射光は偏る：ブルースター角}\par\nobreak
水面やガラスで反射した光は、ある程度偏光しています。マリュスは、1808年、宮殿の窓ガラスで反射した夕日の光を方解石の結晶ごしに見て、このことに気づきました。とくに、入射角がある角 $\theta_{\mathrm B}$ のとき、反射光は、入射面（光線をふくむ紙面）に垂直な向きにゆれる光だけになります。この角を\textbf{ブルースター角}といいます（ブルースター、1815年）。\bkidx{ぶるーすたーかく}{ブルースター角}

\begin{bkfigure}
\begin{tikzpicture}[x=1mm, y=1mm]
  \fill[bkBlue!10] (10,-4) rectangle (100,20);
  \draw[bkInk, line width=0.7pt] (10,20) -- (100,20);
  \node[rxn lab, anchor=south west] at (11,-3.5) {ガラス（屈折率 $n$）};
  \node[rxn lab, anchor=north east] at (99,27) {空気};
  \draw[bkGray, line width=0.4pt, dashed] (50,-2) -- (50,38);
  \draw[rxn ray3] (26.7,35.5) -- (49.5,20.3);
  \draw[rxn ray] (50.4,20.3) -- (73.3,35.5);
  \draw[rxn ray2] (50.2,19.6) -- (64.4,-1.6);
  \draw[bkInk, line width=0.4pt] (51.66,21.11) -- (52.77,19.45) -- (51.11,18.34);
  \draw[bkInk, line width=0.5pt] (50,27) arc (90:146.3:7); \node[rxn lab] at (46.4,30.9) {$\theta_{\mathrm B}$};
  \draw[bkInk, line width=0.5pt] (50,13) arc (-90:-56.3:7); \node[rxn lab] at (52.5,10.3) {$r$};
  \node[rxn lab] at (56.6,18.4) {$90°$};
  % ゆれの印：入射（矢印と点）
  \foreach \px/\py in {34.9/30.1, 43.0/24.7}{\draw[bkInk, line width=0.7pt, {Stealth[length=1.8pt]}-{Stealth[length=1.8pt]}] ({\px-0.555*2.2},{\py-0.832*2.2}) -- ({\px+0.555*2.2},{\py+0.832*2.2});}
  \fill[bkInk] (31.0,32.7) circle (0.7);
  \fill[bkInk] (39.0,27.4) circle (0.7);
  % 反射（点だけ）
  \fill[bkRed!80!black] (58.3,25.6) circle (0.8); \fill[bkRed!80!black] (65.8,30.5) circle (0.8);
  % 屈折（矢印と点）
  \foreach \px/\py in {54.4/13.3, 58.9/6.7}{\draw[bkBlue!80!black, line width=0.7pt, {Stealth[length=1.8pt]}-{Stealth[length=1.8pt]}] ({\px-0.832*2.2},{\py-0.555*2.2}) -- ({\px+0.832*2.2},{\py+0.555*2.2});}
  \fill[bkBlue!80!black] (56.7,10.0) circle (0.7);
  \node[rxn lab, anchor=south] at (24,36) {自然光};
  \node[rxn lab, anchor=south west, align=left] at (73,35.5) {反射光：紙面に垂直なゆれ（点）だけ};
  \node[rxn lab, anchor=north west] at (65,-1.5) {屈折光};
  \node[rxn lab, anchor=west, align=left] at (102,12) {点：紙面に垂直なゆれ\\矢印：紙面の中のゆれ};
\end{tikzpicture}
\bkfigcaption{ブルースター角。反射光と屈折光がちょうど直角になる入射角で、反射光は、紙面に垂直にゆれる光（点の印）だけになる。}
\label{ph:波-13:fig-brewster}
\end{bkfigure}

ブルースター角は、図\ref{ph:波-13:fig-brewster}のように、反射光と屈折光がちょうど直角になる角です。反射角は入射角と同じ $\theta_{\mathrm B}$ なので、屈折角は $r=90°-\theta_{\mathrm B}$ です。屈折の法則 $\sin\theta_{\mathrm B}=n\sin r$（波-12 ③）に入れると、$\sin r=\sin(90°-\theta_{\mathrm B})=\cos\theta_{\mathrm B}$ なので
\[ \sin\theta_{\mathrm B}=n\cos\theta_{\mathrm B}\qquad\text{より}\qquad \tan\theta_{\mathrm B}=n \]
\label{ph:波-13:brewster}となります。ガラス（$n=1.5$）では約 $56°$、水（$n=1.33$）では約 $53°$ です（\nobreak\bkreidairef{11}\nobreak）。

\begin{phyomi}[$\tan\theta_{\mathrm B}=n$]
「屈折率の大きい物質ほど、ブルースター角は大きい（すれすれに近い）」。$n=1$（境がない）なら $45°$ で、反射そのものが起きません。ブルースター角の反射光は、紙面に垂直にゆれる光だけなので、通す向きをそれに直角にした偏光板で、反射を消すことができます。
\end{phyomi}

なぜ直角のとき、紙面の中のゆれが反射しないのでしょう。反射光は、ガラスの表面近くの電子が、屈折光の電場にゆすられて出す光です（⑧と同じしくみ）。ゆれる電子は、ゆれの向きには光を出しません（電-15 で導く結果）。紙面の中のゆれは、屈折光に垂直な向き、つまり反射光の進む向きにゆれているので、その向きへは光が出ないのです。

\begin{bkfigure}
\begin{tikzpicture}
\begin{axis}[bkplot, width=108mm, height=40mm, xmin=0, xmax=90, ymin=0, ymax=1.0, xtick={0,15,30,45,56.3,75,90}, xticklabels={0,15,30,45,\raisebox{-2.2mm}{56.3},75,90}, ytick={0,0.5,1}, xlabel={入射角（°）}, ylabel={反射率}]
  \addplot[bkcurve, domain=0:89, samples=90] {((cos(x)-sqrt(2.25-sin(x)^2))/(cos(x)+sqrt(2.25-sin(x)^2)))^2};
  \addplot[bkcurve2, dashed, domain=0:89, samples=90] {((2.25*cos(x)-sqrt(2.25-sin(x)^2))/(2.25*cos(x)+sqrt(2.25-sin(x)^2)))^2};
  \node[font=\sffamily\footnotesize, anchor=west, align=left, fill=white, inner sep=1.5pt] at (axis cs:4,0.7) {実線：紙面に垂直なゆれ\\破線：紙面の中のゆれ};
\end{axis}
\end{tikzpicture}
\bkfigcaption{ガラス（$n=1.5$）の表面の反射率（電-15 で導く式で計算した）。真上から（$0°$）はどちらも4\,\%。紙面の中のゆれ（破線）は、ブルースター角 $56.3°$ でちょうど0になる。}
\label{ph:波-13:fig-fresnel}
\end{bkfigure}


\begin{phmijika}[偏光サングラスと、カメラのフィルター]
水面や道路の照り返しは、水平にゆれる光を多くふくんでいます（反射面に平行なゆれが反射しやすいから）。偏光サングラスは、通す向きをたてにして、照り返しの光をカットし、水の中の魚や、ぬれた道路を見やすくします。カメラの偏光フィルターも、回して向きを合わせると、ガラスの反射を消したり、青空を濃く写したりできます（青空の光も偏光しているから：⑧）。ミツバチは、空の光の偏光の向きを見て、太陽がかくれていても方角を知ることが知られています。
\end{phmijika}

\bksubsection{7}{干渉計：光の波長でものさしを作る}

\leavevmode\label{ph:波-13:sec7}干渉を使うと、光の波長（約 \qty{0.5}{\micro m}）を目もりにして、長さの変化をはかることができます。その道具が\textbf{干渉計}です。アメリカのマイケルソンが1881年に作った形を、\textbf{マイケルソン干渉計}といいます。\bkidx{まいけるそんかんしょうけい}{マイケルソン干渉計}

\begin{bkfigure}
\begin{tikzpicture}[x=1mm, y=0.8mm]
  \filldraw[fill=bkGray!35, draw=bkInk, line width=0.6pt] (0,16.2) rectangle (16,23.8); \node[rxn lab] at (8,20) {レーザー};
  \draw[bkRed!80!black, line width=1.2pt] (16,20) -- (84,20);
  \draw[bkRed!80!black, line width=1.2pt] (40.6,20.6) -- (40.6,52);
  \draw[bkRed!80!black, line width=1.2pt] (39.4,52) -- (39.4,4);
  \draw[bkRed!80!black, line width=1.2pt, -{Stealth[length=3pt]}] (39.4,8) -- (39.4,4.5);
  \fill[bkBlue!25] (35.8,15.8) -- (36.6,15) -- (44.2,22.6) -- (43.4,23.4) -- cycle; \draw[bkInk, line width=0.5pt] (35.8,15.8) -- (36.6,15) -- (44.2,22.6) -- (43.4,23.4) -- cycle;
  \fill[bkGray!70] (84,12) rectangle (86,28); \node[rxn lab, anchor=west, align=left] at (88,20) {鏡1\\（動かす）};
  \draw[bkInk, line width=0.8pt, {Stealth[length=2.6pt]}-{Stealth[length=2.6pt]}] (82,31) -- (92,31); \node[rxn lab, anchor=south] at (87,31.5) {$\Delta$};
  \fill[bkGray!70] (32,52) rectangle (48,54); \node[rxn lab, anchor=west] at (49,53) {鏡2（止めてある）};
  \fill[bkYellow!40] (32,0) rectangle (48,3); \draw[bkInk, line width=0.5pt] (32,0) rectangle (48,3);
  \node[rxn lab, anchor=west] at (49,1.5) {スクリーン（しまを数える）};
  \node[rxn lab, anchor=north east, align=right] at (34,16) {ハーフミラー\\（半分通し、半分反射）};
  \draw[rxn dim] (42,24.5) -- (84,24.5); \node[rxn lab, anchor=south] at (63,25) {$L_1$};
  \draw[rxn dim] (31,22) -- (31,52); \node[rxn lab, anchor=east] at (30.5,37) {$L_2$};
\end{tikzpicture}
\bkfigcaption{マイケルソン干渉計。レーザーの光をハーフミラーで2つに分け、2つの鏡ではね返った光を、もう一度ハーフミラーで重ねる。鏡1を $\Delta$ だけ動かすと、往復の道のりが $2\Delta$ 変わり、スクリーンのしまが動く。}
\label{ph:波-13:fig-michelson}
\end{bkfigure}

図\ref{ph:波-13:fig-michelson}で、2つの光の経路の差は、往復なので $2(L_1-L_2)$ です。鏡1を $\Delta$ だけ動かすと、経路の差は $2\Delta$ 変わります。経路の差が $\lambda$ 変わるごとに、スクリーンの明るい所と暗い所が1回入れかわり、しまが1本ぶん動きます。だから、しまが $M$ 本動いたとき（③のすきまの本数 $N$ と区別して、動いたしまの本数を $M$ と書きます）
\pheq{2\Delta=M\lambda\qquad\text{つまり}\qquad\Delta=\frac{M\lambda}{2}}{m ＝ m}
\label{ph:波-13:michelson}です。しま1本が、鏡の動き $\lambda/2$ にあたります。赤いレーザー（\qty{633}{nm}）なら、しま1本は約 \qty{0.32}{\micro m} です（\nobreak\bkreidairef{12}\nobreak）。

\begin{phyomi}[$\Delta=M\lambda/2$]
「しまを数えれば、波長の半分を目もりにした長さがわかる」。しまの明るさの変わり方を電気で細かく読めば、1本の何百分の1までわかります。光の波長を目もりにした、ふれずに使えるものさしです。
\end{phyomi}

\begin{phmijika}[レーザーではかる機械の長さと、重力波]
工作機械や半導体を作る装置では、台の位置を、レーザー干渉計で \unit{nm} の細かさではかりながら動かしています。1983年からは、1メートルの長さそのものが、光の速さ（真空で1秒に \qty{299792458}{m}）をもとに決められています。干渉計の究極の例が、重力波の観測装置です。アメリカの LIGO は、長さ \qty{4}{km} の2本の腕をもつ干渉計で、2015年9月、遠くの2つのブラックホールが合体したときに出た重力波が、腕の長さをほんの少し（陽子の大きさの数百分の1）だけのび縮みさせたのを検出しました。日本の神岡には、地下に腕 \qty{3}{km} の KAGRA があります。
\end{phmijika}

\begin{seikaku}
マイケルソンとモーリーは、1887年、干渉計の腕の向きを変えても、しまが動かないことを確かめました。光の速さが、地球の動く向きによって変わらないことを示した実験で、のちの相対性理論の土台の1つになりました（物理シリーズの『原子』（原-11））。
\end{seikaku}

\bksubsection{8}{散乱：ゆすられる電子が、空を青くする}

\leavevmode\label{ph:波-13:sec8}①では、散乱の強さが $1/\lambda^4$ に比例することを、レイリーの法則として使いました。ここでは、なぜ4乗になるのかを、強制振動（波-2）のモデルで導きます。

\Needspace{5\baselineskip}\noindent\textbf{電子は、ばねにつながったおもり}\par\nobreak
空気の分子（窒素・酸素）の中の電子は、原子核に電気の力で引かれて、決まった位置の周りにつなぎとめられています。少しずらすと、もとにもどる力がはたらきます。そこで、電子を「ばねにつながったおもり」とみなします。ばねとおもりの固有角振動数を $\omega_{\mathrm n}$ とします。空気の分子では、$\omega_{\mathrm n}$ は波長でいうと約 \qty{100}{nm} の紫外線の光の角振動数にあたります（目安）。可視光（\qtyrange{380}{780}{nm}）の角振動数 $\omega$ は、その約 $\tfrac14$〜$\tfrac18$ です。

\begin{bkfigure}
\begin{tikzpicture}[x=1mm, y=1mm]
  % 来る光
  \draw[bkRed!80!black, line width=0.9pt, domain=0:40, samples=120, smooth] plot (\x,{14+5*sin(\x*360/13)});
  \draw[bkGray, line width=0.4pt, -{Stealth[length=3pt]}] (0,4) -- (40,4); \node[rxn lab, anchor=north] at (20,3.4) {光の進む向き};
  \foreach \x in {3.25,9.75,16.25,22.75,29.25,35.75}{\pgfmathsetmacro\yy{5*sin(\x*360/13)}\draw[bkRed!60, line width=0.5pt, -{Stealth[length=1.6pt]}] (\x,14) -- (\x,{14+\yy});}
  \node[rxn lab, anchor=south] at (20,21) {光の電場が上下にゆれる};
  % 分子
  \filldraw[fill=bkGray!45, draw=bkInk, line width=0.6pt] (62,6) circle (3.5); \node[rxn lab, anchor=north] at (62,2) {原子核};
  \draw[bkInk, line width=0.6pt, decorate, decoration={coil, aspect=0.4, segment length=1.6mm, amplitude=1.6mm}] (62,9.5) -- (62,19);
  \filldraw[fill=bkBlue!60, draw=bkInk, line width=0.6pt] (62,21) circle (2);
  \node[rxn lab, anchor=west] at (65,33) {電子（質量 $m_{\mathrm e}$）};
  \draw[bkGray, line width=0.4pt] (64.8,32.5) -- (63.2,22.8);
  \draw[bkOrange!85!black, line width=0.9pt, {Stealth[length=2.6pt]}-{Stealth[length=2.6pt]}] (59.5,24.5) -- (59.5,31.5);
  \node[rxn lab, anchor=east] at (58.5,31) {上下にゆすられる};
  % 出る光
  \foreach \a in {-30,0,30}{\draw[bkBlue!80!black, line width=0.8pt, -{Stealth[length=2.6pt]}] (62,21) ++(\a:5) -- ++(\a:8);}
  \foreach \a in {150,180,210}{\draw[bkBlue!80!black, line width=0.8pt, -{Stealth[length=2.6pt]}] (62,21) ++(\a:5) -- ++(\a:7);}
  \node[rxn lab, anchor=west, align=left] at (80,8) {ゆれる電子が、周りへ光を出す\\（ゆれの向き＝上と下には出さない）};
\end{tikzpicture}
\bkfigcaption{散乱のモデル。光の電場が、分子の中の電子を上下にゆする。ゆすられた電子は、同じ振動数の光を周りへ出しなおす。これが散乱した光。}
\label{ph:波-13:fig-denshi}
\end{bkfigure}

\Needspace{5\baselineskip}\noindent\textbf{強制振動の振幅}\par\nobreak
光の電場は、振幅 $E_0$、角振動数 $\omega$ でゆれています。電子は、$E_0$ に比例する大きさ $F_0$ の力で、$F_0\cos\omega t$ のようにゆすられます（$t$ は時刻。②と同じく、④の膜の厚さ $t$ とは別の量です）。ゆれをじゃまする力（減衰）を無視すると、運動方程式は
\[ m_{\mathrm e}\frac{\dd^2x}{\dd t^2}=-m_{\mathrm e}\omega_{\mathrm n}^2\,x+F_0\cos\omega t \]
です。これは強制振動の式（波-2 ④）そのものです。ゆすりに合わせたゆれ（特解）の振幅は、波-2 ④の倍率の式で $\zeta=0$ とおくと
\pheq{A=\frac{F_0/m_{\mathrm e}}{\omega_{\mathrm n}^2-\omega^2}}{m ＝ (N/kg) ÷ (1/s$^2$)}
です（$\omega<\omega_{\mathrm n}$ なので正）。可視光では $\omega$ が $\omega_{\mathrm n}$ よりずっと小さいので、分母はほぼ $\omega_{\mathrm n}^2$ です。つまり、電子のゆれの振幅 $A\approx F_0/(m_{\mathrm e}\omega_{\mathrm n}^2)$ は、光の色（$\omega$）によらず、ほぼ同じです。かたいばねを、ゆっくりゆすっているのと同じで、おもりは力にそのままついていくだけです。

\Needspace{5\baselineskip}\noindent\textbf{出す光は、加速度の2乗に比例する}\par\nobreak
ゆれる電子は、光を出します。電気をもった粒が加速度運動をすると、光（電磁波）を出し、その強さは\textbf{加速度の2乗に比例する}ことが知られています。これは、物理シリーズの『電磁気学』（電-15）で導きます。ここでは、その結果だけを使います。

単振動の加速度は $-\omega^2x$ なので、加速度の振幅は、変位の振幅の $\omega^2$ 倍です（波-1 ②）。だから、散乱した光の強さ $I_{\mathrm s}$ は
\[ I_{\mathrm s}\propto(\omega^2A)^2=\omega^4\Bigl(\frac{F_0/m_{\mathrm e}}{\omega_{\mathrm n}^2-\omega^2}\Bigr)^2 \]
です。$\omega\ll\omega_{\mathrm n}$ では、かっこの中はほぼ一定なので
\[ I_{\mathrm s}\propto\omega^4\propto\frac{1}{\lambda^4} \]
\label{ph:波-13:sanran}となります（$\omega=2\pi c/\lambda$ なので、$\omega^4$ は $1/\lambda^4$ に比例）。①のレイリー散乱の法則が導けました。振幅は色によらないのに、ゆれる速さが速い青い光ほど、電子の加速度が大きく、強く光を出しなおすのです。

\Needspace{10\baselineskip}
\begin{phflow}[空の青さ]
\phflowrow{電子の強制振動\\$m_{\mathrm e}\ddot x=-m_{\mathrm e}\omega_{\mathrm n}^2x+F_0\cos\omega t$}{減衰を無視\\ゆすりに合うゆれ}{$A=\dfrac{F_0/m_{\mathrm e}}{\omega_{\mathrm n}^2-\omega^2}$}
\phflowrow{$I_{\mathrm s}\propto(\omega^2A)^2$\\（電-15 の結果）}{可視光 $\omega\ll\omega_{\mathrm n}$}{$I_{\mathrm s}\propto\dfrac{1}{\lambda^4}$}
\end{phflow}

%% 短い囲みがまたがないように（修3-6、P6-11）
\Needspace{16\baselineskip}
\begin{seikaku}
$\omega$ が $\omega_{\mathrm n}$ に近づくと、分母の $\omega_{\mathrm n}^2-\omega^2$ が小さくなるので、散乱は $1/\lambda^4$ よりも少し強く増えます。$\omega_{\mathrm n}$ を波長 \qty{100}{nm} とすると、\qty{450}{nm} の青と \qty{700}{nm} の赤の比は、$5.9$ ではなく約 $6.2$ になります（グラフ）。実際の空気の測定でも、散乱の強さは $\lambda^{-4}$ より少し急な、$\lambda^{-4.1}$ ほどで変わることが知られています。
\par\vspace{2pt}\centering
\begin{tikzpicture}
\begin{axis}[bkplot, width=100mm, height=40mm, xmin=380, xmax=780, ymin=0, ymax=14, xtick={400,500,600,700}, ytick={0,4,8,12}, xlabel={$\lambda$（nm）}, ylabel={散乱の強さ（\qty{700}{nm} を1）}]
  \addplot[bkcurve] coordinates {(380,12.755) (400,10.240) (420,8.321) (440,6.835) (460,5.669) (480,4.743) (500,4.000) (520,3.398) (540,2.905) (560,2.500) (580,2.163) (600,1.881) (620,1.644) (640,1.443) (660,1.272) (680,1.126) (700,1.000) (720,0.891) (740,0.797) (760,0.715) (780,0.643)};
  \addplot[bkcurve2, dashed, domain=380:780, samples=60] {(700/x)^4};
  \node[font=\sffamily\footnotesize, anchor=west, align=left, fill=white, inner sep=1.5pt] at (axis cs:560,9) {実線：電子のモデル\\破線：$1/\lambda^4$};
\end{axis}
\end{tikzpicture}
\end{seikaku}

\Needspace{5\baselineskip}\noindent\textbf{はてな？の答え}\par\nobreak
\label{ph:波-13:hatena}\bkhatenaref に答えます。空気の分子の中の電子は、太陽の光の電場にゆすられて、強制振動をします。可視光の振動数は、電子の固有振動数よりずっと小さいので、ゆれの振幅はどの色でもほぼ同じです。けれども、青い光は振動数が大きく、電子を速く行ったり来たりさせるので、加速度が大きく、出しなおす光が強くなります。その強さは振動数の4乗に比例し、青（\qty{450}{nm}）は赤（\qty{700}{nm}）の約 5.9 倍です。だから、太陽のない方向から目に届く空の光は、青を多くふくみ、空は青く見えます。

夕方は、太陽の光が空気の中を長く通るので（図\ref{ph:波-13:fig-yuuyake}）、途中で青い光が次々に散乱されて抜けていきます。光が空気の層を通って、どれだけ減るかは、「同じ厚さの空気を通るごとに、同じ割合で減る」と考えて計算できます（\nobreak\bkreidairef{13}\nobreak）。真上から来る光に比べて、地平線の近くから来る光は、空気の中を約38倍の道のりで通ります。そのあいだに、青い光はほとんど残らず、赤い光が約4分の1残ります。夕日が赤いのは、このためです。

\begin{phmijika}[空はなぜ紫ではない？　雲はなぜ白い？]
散乱は、紫（\qty{400}{nm}）のほうが青よりさらに強いのに、空は紫に見えません。太陽の光には、もともと紫の光が青より少ないこと（熱-13 ③の太陽の光の分布）と、人の目が紫の光を感じにくいことが主な理由です。散らされた光は、紫・青・緑をふくみ、目にはそれがまざった空色に見えます。

雲の水の粒は直径 数〜数十 \unit{\micro m} で、光の波長よりずっと大きいので、電子1つ1つのモデルは使えません。大きな粒の散乱は、波長にほとんどよらないので、どの色も同じように散らされ、雲は白く見えます。また、太陽から $90°$ の向きの青空の光は、偏光しています（ゆすられる電子は、横から見ると1つの向きにしかゆれて見えないから。⑥）。偏光フィルターで青空が濃く写るのは、このためです。
\end{phmijika}

%% ---------------- たいせつ・ことば・記号 ----------------
\begin{taisetsu}
\begin{itemize}[leftmargin=1.2\zw, itemsep=2pt, topsep=0pt, label={\color{bkNavy}●}]\raggedright
\item 光の明るい所・暗い所は、経路（光路長）の差で決まる。差が波長の整数倍で強め合い、整数＋$\tfrac12$ 倍で弱め合う。物質の中の長さは $n$ 倍して数える。
\item ヤングの実験 $\Delta x=L\lambda/d$。回折格子 $d\sin\theta=m\lambda$（本数が多いほど線が細い）。
\item 薄い膜：光路長の差 $2nt\cos r$ と、屈折率の大きい物質との境での反射の入れかわり（半波長）を、1つずつ数える。くさび $\Delta x=L\lambda/(2H)$、ニュートンリング $\rho_m^{\,2}=m\lambda R$。
\item 円い穴の回折：はじめて暗くなる向きは $\sin\theta=1.22\lambda/D$。見分けられる最小の角 $\theta_{\min}\approx1.22\lambda/D$。
\item 偏光：自然光は偏光板で $\tfrac12$。マリュスの法則 $I=I_0\cos^2\theta$。ブルースター角 $\tan\theta_{\mathrm B}=n$。干渉計：しま1本が鏡の動き $\lambda/2$。
\item 散乱：ゆすられる電子（強制振動）が出す光は $\omega^4\propto1/\lambda^4$。空は青く、夕日は赤い。
\end{itemize}
\end{taisetsu}

\begin{kotoba}\small
\bkword{干渉じま}{2つ以上の光が重なってできる、明るい所と暗い所の並び。}
\bkword{次数}{明るい線・暗い線が、まん中から何番めかを表す整数 $m$。}
\bkword{回折格子・分光}{細い溝を等間かくに並べた板。光を波長ごとに分けること。}
\bkword{光路長}{屈折率×物質の中の長さ。真空なら何 m ぶんの波が入っているか。}
\bkword{反射防止膜}{2つの反射光を弱め合わせて、反射を減らす薄い膜。}
\bkword{ニュートンリング}{レンズと板のあいだの空気の層でできる、同心円の干渉じま。}
\bkword{分解能}{近い2つの点を見分けられる細かさ。円い穴では $1.22\lambda/D$。}
\bkword{偏光}{ゆれの向きがそろった光。光が横波であることの証拠。}
\bkword{ブルースター角}{反射光が完全に偏光する入射角。$\tan\theta_{\mathrm B}=n$。}
\bkword{干渉計}{光を分けて重ね、しまの動きから長さの変化をはかる道具。}
\bkword{レイリー散乱}{波長よりずっと小さな粒による散乱。強さは $1/\lambda^4$ に比例。}
\end{kotoba}
\bkkigocards{\bkkigocard{$m$}{エム}{干渉の次数（整数）。明るい線・暗い線のまん中からの番号。波-13 では、電子の質量を $m_{\mathrm e}$ と書いて区別する}{order of interference}{波-13}}{\bkkigocard{$\theta_{\mathrm B}$}{シータ・ビー}{ブルースター角。反射光が完全に偏光する入射角。$\tan\theta_{\mathrm B}=n$}{Brewster angle}{波-13}}

\begin{senmonsho}[回折限界と薄膜干渉]
\begin{bkhonbun}
口径 $D$ の円形開口による点像の強度分布はエアリーパターンとなり、第1暗環の角半径は $\theta\simeq1.22\lambda/D$ である。これを光学系の角分解能（レイリーの分解能）という。また、膜厚 $d$、屈折率 $n$ の薄膜における反射光の干渉条件は、位相の反転を考慮して $2nd\cos\phi=(m+1/2)\lambda$ と表される。
\end{bkhonbun}
\bkyaku{口径＝穴（レンズの口）の直径。点像＝点の像。第1暗環＝はじめの暗い輪。$\simeq$ は「ほぼ等しい」。膜厚＝膜の厚さ（この本の $t$）、$\phi$ は屈折角（この本の $r$）。「位相の反転」は、この本の「山と谷の入れかわり」。$2nd\cos\phi=(m+1/2)\lambda$ は、反射光が強め合う条件（④の式）。}
\end{senmonsho}

%% ---------------- 例題 ----------------
\bksection{reidai}{例題でたしかめよう}

\noindent 光の速さは $c=3.00\times10^8$\,m/s、角は小さいとき $\sin\theta\approx\tan\theta\approx\theta$（rad）とします。長さの単位は、計算の前に m にそろえます（\qty{1}{nm}＝$10^{-9}$\,m、\qty{1}{\micro m}＝$10^{-6}$\,m、\qty{1}{mm}＝$10^{-3}$\,m）。

\begin{phreidai}[空の青と、赤い光の散乱を比べる]
\label{ph:波-13:ex1}空気の分子による散乱で、青い光（\qty{450}{nm}）は、赤い光（\qty{700}{nm}）の何倍強く散らされますか。緑の光（\qty{550}{nm}）は、赤の何倍ですか。
\phbamen
\bkhatenaref の場面です。太陽の白い光が空気の中を通るとき、空気の分子に当たって、四方へ散らされます（図\ref{ph:波-13:fig-yuuyake}）。
\phmushi
太陽の光にふくまれる色ごとの強さの違い、目の感じ方の違い、ちりや水の粒による散乱を無視します（空気の分子だけを考える）。
\phzu
青・緑・赤の波長を $\lambda_{\text{青}}=\qty{450}{nm}$、$\lambda_{\text{緑}}=\qty{550}{nm}$、$\lambda_{\text{赤}}=\qty{700}{nm}$ とする。
\phriyu
空気の分子は光の波長よりずっと小さい（約 \qty{0.3}{nm}）ので、レイリー散乱で、強さは $1/\lambda^4$ に比例する（①・⑧）。比べるだけなので、比例の定数はいらない。
\phtateru
\[ \frac{\text{青の散乱}}{\text{赤の散乱}}=\frac{1/\lambda_{\text{青}}^4}{1/\lambda_{\text{赤}}^4}=\Bigl(\frac{\lambda_{\text{赤}}}{\lambda_{\text{青}}}\Bigr)^4 \]
\phtoku
$700\div450\approx1.556$。4乗は、2乗を2回して $1.556^2\approx2.420$、$2.420^2\approx5.86$。緑は $700\div550\approx1.273$、$1.273^2\approx1.620$、$1.620^2\approx2.62$。
\phtashikame
単位：nm÷nm で単位なし。大きさ：波長の短い順に、青＞緑＞赤と強く散らされ、①のグラフと合う。極端な場合：同じ波長どうしなら比は1。
\kotae 青は赤の約5.9倍、緑は赤の約2.6倍
\end{phreidai}
\begin{jibun}
紫の光（\qty{400}{nm}）は、赤い光（\qty{700}{nm}）の何倍強く散らされますか。
\bkjans{$(700\div400)^4=1.75^4\approx9.4$。約9.4倍}
\end{jibun}

\begin{phreidai}[ヤングの実験のしまの間かく]
\label{ph:波-13:ex2}波長 \qty{633}{nm} の赤いレーザーの光を、間かく \qty{0.20}{mm} の2つのすきまに当て、\qty{1.0}{m} 離れたスクリーンで干渉じまを見ます。(1) 明るい線の間かくを求めましょう。(2) まん中から3本め（$m=3$）の明るい線は、まん中から何 mm の所にありますか。
\phbamen
図\ref{ph:波-13:fig-young}の実験です。
\phmushi
すきまのはば、空気の屈折率（1とする）、スクリーンの上での光の弱まりを無視します。
\phzu
$d=\qty{0.20}{mm}=2.0\times10^{-4}$\,m、$L=\qty{1.0}{m}$、$\lambda=633\times10^{-9}$\,m。明るい線の間かく $\Delta x$、3本めの位置 $x_3$。
\phriyu
$L$ は $d$ の5000倍で十分遠く、角も小さいので、②の $\Delta x=L\lambda/d$ が使える。明るい線は等間かくに並ぶので、3本めは $3\Delta x$。
\phtateru
\[ \Delta x=\frac{L\lambda}{d},\qquad x_3=3\,\Delta x \]
\phtoku
\[ \Delta x=\frac{1.0\times633\times10^{-9}}{2.0\times10^{-4}}=3.165\times10^{-3}\ \mathrm{m}\approx\qty{3.2}{mm} \]
$x_3=3\times\qty{3.165}{mm}\approx\qty{9.5}{mm}$。
\phtashikame
単位：m×m÷m＝m。大きさ：mm ほどの間かくで、目で数えられる。角は $9.5\times10^{-3}\div1.0\approx0.0095$\,rad と小さく、$\sin\theta\approx\tan\theta$ の近似のずれは0.01\,\%以下。極端な場合：$d$ を大きくすると $\Delta x$ は小さくなり、しまは見えなくなる。
\kotae (1) 約 \qty{3.2}{mm}　(2) 約 \qty{9.5}{mm}
\end{phreidai}
\begin{jibun}
すきまの間かく \qty{0.25}{mm}、スクリーンまで \qty{2.0}{m} の装置に、色のわからないレーザーを当てると、明るい線の間かくは \qty{4.3}{mm} でした。光の波長を求めましょう。
\bkjans{$\lambda=\Delta x\,d/L=4.3\times10^{-3}\times0.25\times10^{-3}\div2.0$ $\approx5.4\times10^{-7}$\,m。約 \qty{540}{nm}（緑）}
\end{jibun}

\begin{phreidai}[回折格子で白い光を分ける]
\label{ph:波-13:ex3}1 mm あたり1000本の溝のある回折格子に、白い光を垂直に当てます。(1) 波長 \qty{550}{nm} の光の1次（$m=1$）の明るい線の向きを求めましょう。(2) 1次の虹の帯（\qtyrange{400}{700}{nm}）は、何度から何度までに広がりますか。(3) \qty{550}{nm} の2次の明るい線はできますか。
\phbamen
回折格子を使った分光器です（図\ref{ph:波-13:fig-bunkou}）。
\phmushi
溝の形と深さ（明るさだけに効く）、格子の厚さを無視します。明るい線は、遠くで（またはレンズで集めて）見るとします。
\phzu
溝の間かく $d=\tfrac{1}{1000}$\,mm $=1.0\times10^{-6}$\,m。角 $\theta$、次数 $m$。
\phriyu
$N$ 本のすきまが同じ間かくで並んでいるので、③の $d\sin\theta=m\lambda$。
\phtateru
\[ \sin\theta=\frac{m\lambda}{d} \]
\phtoku
(1) $\sin\theta=550\times10^{-9}\div(1.0\times10^{-6})=0.55$。電卓の $\sin^{-1}$ で $\theta\approx33.4°$。\par
(2) \qty{400}{nm}：$\sin\theta=0.40$ で $\theta\approx23.6°$。\qty{700}{nm}：$\sin\theta=0.70$ で $\theta\approx44.4°$。\par
(3) $m=2$ では $\sin\theta=2\times0.55=1.10$。$\sin\theta$ は1より大きくなれないので、この向きはない。
\phtashikame
$\sin\theta$ は単位なし（m÷m）。大きさ：$d$ が波長の2倍ほどしかないので、大きく曲がる。紫が内側、赤が外側で、図\ref{ph:波-13:fig-bunkou}と同じ並び。極端な場合：$d$ を大きく（本数を少なく）すると、角は小さくなり、高い次数まで現れる。
\kotae (1) 約 $33°$　(2) 約 $24°$ から $44°$ まで　(3) できない（$\sin\theta>1$）
\end{phreidai}
\begin{jibun}
1 mm あたり600本の回折格子に、ナトリウムランプの黄色い光を当てると、1次の明るい線が $20.7°$ の向きにできました。光の波長を求めましょう。
\bkjans{$d=\tfrac{1}{600}$\,mm $\approx1.667\times10^{-6}$\,m。$\lambda=d\sin20.7°$ $\approx1.667\times10^{-6}\times0.3535$ $\approx5.89\times10^{-7}$\,m。約 \qty{590}{nm}}
\end{jibun}

\begin{bkfigure}
\begin{tikzpicture}[x=1.6mm, y=1.3mm]
  \fill[bkGray!30] (-1.2,-18) rectangle (0,18); \node[rxn lab, anchor=north] at (-0.6,-18.6) {CD};
  \fill[bkYellow!25] (30,-18) rectangle (30.8,-0.6); \fill[bkYellow!25] (30,0.6) rectangle (30.8,18);
  \draw[bkInk, line width=0.4pt] (30,-18) rectangle (30.8,18);
  \node[rxn lab, anchor=north] at (30.4,-18.6) {紙（穴あき）};
  \filldraw[fill=bkGray!40, draw=bkInk, line width=0.5pt] (34,-1.5) rectangle (42,1.5);
  \node[rxn lab, anchor=north] at (38,-2.2) {レーザー};
  \draw[bkRed!80!black, line width=1pt] (34,0) -- (0,0);
  \draw[bkRed!70, line width=0.6pt] (0,0) -- (30,13.35); \draw[bkRed!70, line width=0.6pt] (0,0) -- (30,-13.35);
  \draw[bkRed!50, line width=0.6pt, dashed] (0,0) -- (12.95,18); \draw[bkRed!50, line width=0.6pt, dashed] (0,0) -- (12.95,-18);
  \fill[bkRed] (30,13.35) circle (0.6); \fill[bkRed] (30,-13.35) circle (0.6);
  \draw[bkInk, line width=0.5pt] (7,0) arc (0:24:7); \node[rxn lab] at (9.6,1.9) {$\theta$};
  \draw[rxn dim] (33,0.8) -- (33,13.35); \node[rxn lab, anchor=west] at (33.4,8) {$y$};
  \draw[rxn dim] (0,-15.5) -- (30,-15.5); \node[rxn lab, fill=white, inner sep=0.5pt] at (15,-15.5) {\qty{0.30}{m}};
  \node[rxn lab, anchor=south west] at (31.2,14.2) {1次の点};
  \node[rxn lab, anchor=west] at (13.6,17) {2次（破線）};
\end{tikzpicture}
\bkfigcaption{CD にレーザーを垂直に当てる（上から見たところ）。トラックで回折した光が、紙の上に、穴の両側に対称な光の点を作る。}
\label{ph:波-13:fig-cd}
\end{bkfigure}

\begin{phreidai}[CD に赤いレーザーを当てる]
\label{ph:波-13:ex4}赤いレーザーポインター（\qty{650}{nm}）の光を、CD の記録面に垂直に当てます。CD の正面 \qty{0.30}{m} の所に、CD と平行に白い紙を立て、紙の小さな穴からレーザーを通します（図\ref{ph:波-13:fig-cd}）。トラックの間かくは \qty{1.6}{\micro m} です。紙に映る1次の光の点の向きと、穴からの距離を求めましょう。2次の点はどうですか。（レーザーの光を目に入れないこと。）
\phbamen
CD の細かいトラックの並びが、光をはね返す回折格子になります。やってみようの③を、1つの色の光で数にします。
\phmushi
トラックがうずまき状でわずかに曲がっていること（光が当たる小さな範囲では、まっすぐな線の並びとみなす）、CD の透明な層での屈折を無視します。
\phzu
$d=1.6\times10^{-6}$\,m、$\lambda=650\times10^{-9}$\,m、紙までの距離 \qty{0.30}{m}、穴から点までの距離 $y$。
\phriyu
となりのトラックではね返った光の経路の差も $d\sin\theta$ なので、反射でも③の $d\sin\theta=m\lambda$ が使える。紙の上の距離は、直角三角形で $y=0.30\tan\theta$。
\phtateru
\[ \sin\theta=\frac{m\lambda}{d},\qquad y=0.30\tan\theta \]
\phtoku
$m=1$：$\sin\theta=0.65\div1.6\approx0.406$（\unit{\micro m} どうしでわった）、$\theta\approx24.0°$。\par
$y=0.30\times\tan24.0°\approx0.30\times0.445\approx\qty{0.13}{m}$。\par
$m=2$：$\sin\theta\approx0.813$、$\theta\approx54.3°$、$y\approx0.30\times1.39\approx\qty{0.42}{m}$。\par
$m=3$：$\sin\theta\approx1.22>1$ なので、できない。
\phtashikame
単位：$y$ は m。大きさ：紙の上で十数 cm 離れた所に点が見え、家でも確かめられる。左右（上下）対称に2つずつできる。極端な場合：トラックの間かくがせまいほど、点は外へ離れる（DVD）。
\kotae 1次は約 $24°$、穴から約 \qty{13}{cm}。2次は約 $54°$、約 \qty{42}{cm}（紙が大きければ見える）
\end{phreidai}
\begin{jibun}
トラックの間かくが \qty{0.74}{\micro m} の DVD に、同じように \qty{650}{nm} の光を当てます。1次の点の向きと、穴からの距離（紙まで \qty{0.30}{m}）を求めましょう。2次の点はできますか。
\bkjans{$\sin\theta=0.65\div0.74\approx0.878$、$\theta\approx61.4°$。$y=0.30\times\tan61.4°$ $\approx\qty{0.55}{m}$。2次は $\sin\theta\approx1.76>1$ でできない}
\end{jibun}

\begin{phreidai}[シャボン膜の色]
\label{ph:波-13:ex5}厚さ \qty{300}{nm} のシャボン膜（屈折率1.33）を、ほぼ垂直に、白い光で見ます。(1) 強く反射する可視光の波長を求めましょう。(2) 膜が \qty{10}{nm} まで薄くなると、何色に見えますか。
\phbamen
シャボン玉やシャボン膜の色です（図\ref{ph:波-13:fig-maku}）。
\phmushi
膜の厚さのむら、光のななめの入り方（垂直とする）、3回以上反射する光を無視します。
\phzu
$n=1.33$、$t=\qty{300}{nm}$、強め合う波長 $\lambda$、整数 $m$。
\phriyu
表の面（空気→膜）で入れかわり、裏の面（膜→空気）ではそのまま。入れかわりは1つなので、④の「強め合う：$2nt=(m+\tfrac12)\lambda$」を使う。
\phtateru
\[ \lambda=\frac{2nt}{m+\frac12}\qquad(m=0,\ 1,\ 2,\ \cdots) \]
\phtoku
(1) $2nt=2\times1.33\times300=\qty{798}{nm}$。$m=0$：$798\div0.5=\qty{1596}{nm}$（赤外線）。$m=1$：$798\div1.5=\qty{532}{nm}$（緑）。$m=2$：$798\div2.5\approx\qty{319}{nm}$（紫外線）。可視光は \qty{532}{nm} だけ。（弱め合うのは $2nt=m\lambda$ の \qty{399}{nm}（紫のはし）と \qty{798}{nm}（赤外線）。）\par
(2) $2nt=2\times1.33\times10=\qty{26.6}{nm}$ は、可視光の波長（\qty{380}{nm} 以上）の $\tfrac{1}{14}$ 以下。経路の差はほぼ0で、反射の入れかわりの半波長だけが残るので、どの色も弱め合う。
\phtashikame
単位：nm。$m$ の候補を全部調べ、可視光に入るものだけを選んだ。極端な場合：$t\to0$ では、どの色も弱め合って黒くなり、④の「黒い膜」と合う。
\kotae (1) 約 \qty{530}{nm}（緑に見える）　(2) どの色も弱め合い、黒く見える
\end{phreidai}
\begin{jibun}
同じシャボン膜が \qty{200}{nm} の厚さになりました。ほぼ垂直に見るとき、弱め合う可視光の波長を求めましょう。
\bkjans{$2nt=2\times1.33\times200=\qty{532}{nm}$。弱め合うのは $2nt=m\lambda$ の $m=1$ で約 \qty{532}{nm}（緑が消え、赤紫っぽく見える）}
\end{jibun}

\begin{phreidai}[カメラのレンズの反射防止膜]
\label{ph:波-13:ex6}ガラス（屈折率1.52）のレンズに、フッ化マグネシウム（屈折率1.38）の膜をつけて、緑の光（\qty{550}{nm}）の反射を弱めます。いちばん薄い膜の厚さを求めましょう。
\phbamen
\noindent\begin{minipage}[c]{0.56\linewidth}\raggedright
カメラやめがねのレンズです。膜の表と裏で反射した2つの光を、弱め合わせます。
\end{minipage}\hfill
\begin{minipage}[c]{0.4\linewidth}\centering
\begin{tikzpicture}[x=1mm, y=1mm]
  \fill[bkPurple!12] (0,8) rectangle (40,13); \fill[bkBlue!14] (0,0) rectangle (40,8);
  \draw[bkInk, line width=0.6pt] (0,13) -- (40,13); \draw[bkInk, line width=0.6pt] (0,8) -- (40,8);
  \node[rxn lab, anchor=west] at (24,10.5) {膜 1.38};
  \node[rxn lab, anchor=west] at (24,4) {ガラス 1.52};
  \node[rxn lab, anchor=west] at (24,17) {空気 1.00};
  \draw[rxn ray3] (2,24) -- (7.6,13.4);
  \draw[rxn ray] (8.2,13.4) -- (12,21);
  \draw[bkBlue!70!black, line width=0.8pt] (8,13) -- (10.5,8.2) -- (13,13);
  \draw[rxn ray2] (13,13.4) -- (16.8,21);
  \node[rxn lab, anchor=south, align=center] at (14,21.5) {どちらも\\入れかわる};
\end{tikzpicture}
\end{minipage}
\phmushi
2回以上反射する光、光のななめの入り方（垂直とする）、膜の中での光の吸収を無視します。
\phzu
膜の屈折率 $n=1.38$、厚さ $t$、$\lambda=\qty{550}{nm}$。
\phriyu
表の面（$1.00\to1.38$）でも、裏の面（$1.38\to1.52$）でも、屈折率の大きい物質との境で反射するので、どちらも入れかわる。2つの入れかわりは打ち消し合うので、光路長の差 $2nt$ だけで決まり、弱め合う条件は $2nt=(m+\tfrac12)\lambda$。いちばん薄いのは $m=0$。
\phtateru
\[ 2nt=\frac12\lambda\qquad\text{より}\qquad t=\frac{\lambda}{4n} \]
\phtoku
\[ t=\frac{550}{4\times1.38}=\frac{550}{5.52}\approx\qty{99.6}{nm} \]
\phtashikame
単位：nm。大きさ：膜の中の波長は $550\div1.38\approx\qty{399}{nm}$ で、その4分の1。往復すると半波長ずれて、ちょうど打ち消し合う。次に弱め合う厚さは $3\lambda/(4n)\approx\qty{299}{nm}$。
\kotae 約 \qty{100}{nm}
\end{phreidai}
\begin{jibun}
同じフッ化マグネシウムの膜で、だいだい色の光（\qty{600}{nm}）の反射をいちばん弱めるには、いちばん薄くて何 nm にしますか。
\bkjans{$t=600\div(4\times1.38)\approx\qty{109}{nm}$（約 \qty{110}{nm}）}
\end{jibun}

\begin{phreidai}[髪の毛の太さを光ではかる]
\label{ph:波-13:ex7}長さ \qty{10}{cm} の2枚の平らなガラス板の一方のはしに、髪の毛を1本はさみます。ナトリウムランプの光（\qty{589}{nm}）を上から当てて見ると、暗い線が \qty{0.50}{mm} の間かくで並びました。髪の毛の太さを求めましょう。
\phbamen
図\ref{ph:波-13:fig-kusabi}のくさび形の空気の層です。
\phmushi
ガラス板のたわみ、光のななめの入り方、ガラス板の外側の面での反射を無視します。
\phzu
板の長さ $L=0.10$\,m、暗い線の間かく $\Delta x=0.50\times10^{-3}$\,m、$\lambda=589\times10^{-9}$\,m、髪の毛の太さ $H$。
\phriyu
入れかわりは1つ（下の板の上の面）なので、暗い線は $2t=m\lambda$ の所にできる。④の $\Delta x=L\lambda/(2H)$ を $H$ について解く。
\phtateru
\[ \Delta x=\frac{L\lambda}{2H}\qquad\text{より}\qquad H=\frac{L\lambda}{2\,\Delta x} \]
\phtoku
\[ H=\frac{0.10\times589\times10^{-9}}{2\times0.50\times10^{-3}}=5.89\times10^{-5}\ \mathrm{m}\approx\qty{59}{\micro m} \]
\phtashikame
単位：m×m÷m＝m。\\
別の見方：\qty{10}{cm} の中の暗い線の数は
\[ 0.10\div(0.50\times10^{-3})=200\ \text{本} \]
で、1本ごとに厚さが $\lambda/2$ ずつ増えるので、$200\times\qty{294.5}{nm}\approx\qty{59}{\micro m}$。同じ答えになった。\\
大きさ：髪の毛は \qtyrange{0.05}{0.1}{mm} ほどで、合っている。
\kotae 約 \qty{59}{\micro m}（約 \qty{0.06}{mm}）
\end{phreidai}
\begin{jibun}
長さ \qty{5.0}{cm} の2枚のガラス板のはしに、厚さ \qty{74}{\micro m} の紙をはさみ、\qty{589}{nm} の光を上から当てます。暗い線の間かくを求めましょう。
\bkjans{$\Delta x=L\lambda/(2H)$ $=0.050\times589\times10^{-9}$ $\div(2\times74\times10^{-6})$ $\approx2.0\times10^{-4}$\,m。約 \qty{0.20}{mm}}
\end{jibun}

\begin{phreidai}[ニュートンリングでレンズの丸みをはかる]
\label{ph:波-13:ex8}平凸レンズを平らなガラス板にのせ、\qty{589}{nm} の光で上から見ると、中心から数えて10番めの暗い輪（$m=10$）の半径が \qty{2.4}{mm} でした。レンズの球面の半径 $R$ を求めましょう。
\phbamen
図\ref{ph:波-13:fig-newton}のニュートンリングです。レンズの工場で、球面の丸みを調べる方法です。
\phmushi
レンズの重さによるつぶれ（中心で、ちょうどふれ合っているとする）、光のななめの入り方、ちりのはさまりを無視します。
\phzu
$\rho_{10}=2.4\times10^{-3}$\,m、$m=10$、$\lambda=589\times10^{-9}$\,m、球面の半径 $R$。
\phriyu
空気の層の厚さは $t\approx\rho^2/(2R)$、入れかわりは1つで、暗い輪は $2t=m\lambda$。④の $\rho_m^{\,2}=m\lambda R$ を $R$ について解く。
\phtateru
\[ R=\frac{\rho_m^{\,2}}{m\lambda} \]
\phtoku
\[ R=\frac{(2.4\times10^{-3})^2}{10\times589\times10^{-9}}=\frac{5.76\times10^{-6}}{5.89\times10^{-6}}\approx\qty{0.98}{m} \]
\phtashikame
単位：m$^2$÷m＝m。近似の確かめ：輪の所の空気の厚さは $t=\rho^2/(2R)\approx\qty{2.9}{\micro m}$ で、$R$ の30万分の1。$t^2$ を無視してよい。$2t\approx\qty{5.9}{\micro m}$ は $10\lambda$ で、10番めの輪と合う。
\kotae 約 \qty{0.98}{m}
\end{phreidai}
\begin{jibun}
球面の半径 \qty{1.0}{m} の平凸レンズで、\qty{589}{nm} の光のニュートンリングを見ます。4番めの暗い輪の半径を求めましょう。
\bkjans{$\rho_4=\sqrt{4\times589\times10^{-9}\times1.0}$ $\approx1.5\times10^{-3}$\,m。約 \qty{1.5}{mm}}
\end{jibun}

\begin{phreidai}[目で見分けられる限界]
\label{ph:波-13:ex9}明るい所で、ひとみの直径を \qty{3.0}{mm}、光の波長を \qty{550}{nm} とします。(1) 目が見分けられる最小の角を求めましょう。(2) \qty{1.0}{mm} 離れた2つの点を見分けられるのは、目から何 m までですか。
\phbamen
\noindent\begin{minipage}[c]{0.56\linewidth}\raggedright
視力検査の輪の切れ目のように、近い2つの点を見分けます。ひとみは円い穴です。
\end{minipage}\hfill
\begin{minipage}[c]{0.4\linewidth}\centering
\begin{tikzpicture}[x=1mm, y=1mm]
  \fill[bkInk] (0,2) circle (0.8); \fill[bkInk] (0,-2) circle (0.8);
  \draw[bkInk, line width=0.5pt] (0,2) -- (26,0) -- (0,-2);
  \rxneye{(29,0)}
  \draw[bkInk, line width=0.5pt] (18,0.62) arc (175.6:184.4:8); \node[rxn lab, anchor=south] at (18,2.0) {$\theta$};
  \draw[rxn dim] (-2.5,-2) -- (-2.5,2); \node[rxn lab, anchor=east] at (-3,0) {\qty{1.0}{mm}};
  \draw[rxn dim] (0,-5) -- (26,-5); \node[rxn lab, anchor=north] at (13,-5.5) {距離 $L$};
\end{tikzpicture}
\end{minipage}
\phmushi
目の中の液体で波長が少し短くなること、目のレンズのゆがみ（収差）、網膜の細胞の大きさを無視します。
\phzu
$D=3.0\times10^{-3}$\,m、$\lambda=550\times10^{-9}$\,m、見分けられる最小の角 $\theta_{\min}$、距離 $L$。
\phriyu
ひとみは円い穴なので、⑤の $\theta_{\min}\approx1.22\lambda/D$。2つの点の間かくは、角が小さいので $L\theta$（弧の長さ $r\theta$ と同じ、4-2）。
\phtateru
\[ \theta_{\min}=1.22\,\frac{\lambda}{D},\qquad L\,\theta_{\min}=1.0\times10^{-3} \]
\phtoku
\[ \theta_{\min}=\frac{1.22\times550\times10^{-9}}{3.0\times10^{-3}}\approx2.24\times10^{-4}\ \mathrm{rad} \]
度に直すと $2.24\times10^{-4}\times\dfrac{180}{\pi}\approx0.0128°$、分に直すと $0.0128\times60\approx0.77$ 分。
\[ L=\frac{1.0\times10^{-3}}{2.24\times10^{-4}}\approx\qty{4.5}{m} \]
\phtashikame
単位：m÷m で rad。大きさ：視力1.0の人が見分ける角は約1分なので、目は回折の限界に近い所まで見えている。極端な場合：暗い所でひとみが広がると、回折の限界は小さくなる。
\kotae (1) 約 $2.2\times10^{-4}$\,rad（約0.8分）　(2) 約 \qty{4.5}{m}
\end{phreidai}
\begin{jibun}
口の直径 \qty{0.10}{m} の望遠レンズで、\qty{550}{nm} の光を使って \qty{1.0}{km} 先を写します。見分けられる2つの点の、最小の間かくを求めましょう。
\bkjans{$\theta_{\min}=1.22\times550\times10^{-9}\div0.10$ $\approx6.7\times10^{-6}$\,rad。$1000\times6.7\times10^{-6}$ $\approx6.7\times10^{-3}$\,m。約 \qty{6.7}{mm}}
\end{jibun}

\begin{phreidai}[偏光板を重ねる]
\label{ph:波-13:ex10}(1) 自然光（強さ $I_0$）を偏光板1に通したあと、通す向きが $60°$ 違う偏光板2に通します。偏光板2を出た光の強さは、$I_0$ の何\,\%ですか。(2) 通す向きが直角の2枚の偏光板の間に、通す向きが $45°$ の偏光板をはさむと、出てくる光は $I_0$ の何\,\%ですか。
\phbamen
図\ref{ph:波-13:fig-henkou}のように、偏光サングラスを重ねて回す場面です。
\phmushi
偏光板が、通す向きの光も少し吸うこと（理想の偏光板とする）、偏光板の表面での反射を無視します。
\phzu
はじめの強さ $I_0$。1枚めを出た光 $I_1$、2枚め $I_2$、3枚め $I_3$。
\phriyu
自然光は、1枚めでどの向きの $\cos^2$ も平均されて $\tfrac12$。その後は、ゆれの向きがそろった光なので、マリュスの法則 $I=I_0\cos^2\theta$ を1枚ごとに使う（⑥）。
\phtateru
(1) $I_1=\tfrac12I_0$、$I_2=I_1\cos^260°$。(2) $I_1=\tfrac12I_0$、$I_2=I_1\cos^245°$、$I_3=I_2\cos^245°$（2枚めと3枚めの向きの差も $45°$）。
\phtoku
(1) $I_2=\tfrac12I_0\times0.25=0.125\,I_0$。\par
(2) $I_3=\tfrac12I_0\times0.5\times0.5=0.125\,I_0$。
\phtashikame
割合はどれも1以下で、通るたびに減っている。極端な場合：(1) で $60°$ を $0°$ にすると $\tfrac12$、$90°$ にすると0。(2) では、間の1枚をぬくと直角の2枚になり、光はまったく通らない。1枚\textbf{加える}と明るくなるのは、間の偏光板が、光のゆれの向きを少しずつ回して、次の偏光板を通れる成分を作るからです。
\kotae (1) 12.5\,\%（$\tfrac18$）　(2) 12.5\,\%（間の1枚がないと0）
\end{phreidai}
\begin{jibun}
(1) 自然光を、通す向きが $30°$ 違う2枚の偏光板に通すと、強さは $I_0$ の何\,\%になりますか。(2) 通す向きが直角の2枚の間に、1枚めから $30°$ の向きの偏光板をはさむと、何\,\%ですか。
\bkjans{(1) $\tfrac12\times\cos^230°=\tfrac12\times0.75=0.375$。37.5\,\%。(2) $\tfrac12\times\cos^230°\times\cos^260°$ $=0.5\times0.75\times0.25\approx0.094$。約9.4\,\%}
\end{jibun}

\begin{phreidai}[窓ガラスのブルースター角]
\label{ph:波-13:ex11}ガラス（屈折率1.50）の表面で反射した光が、完全に偏光するときの入射角と、そのときの屈折角を求めましょう。
\phbamen
図\ref{ph:波-13:fig-brewster}の場面です。窓ガラスの照り返しを、偏光フィルターで消すときに使います。
\phmushi
ガラスの裏の面での反射、ガラスによる光の吸収を無視します。
\phzu
屈折率 $n=1.50$、ブルースター角 $\theta_{\mathrm B}$、屈折角 $r$。
\phriyu
反射光と屈折光が直角になる角で、⑥の $\tan\theta_{\mathrm B}=n$ が成り立つ。屈折角は $r=90°-\theta_{\mathrm B}$。
\phtateru
\[ \theta_{\mathrm B}=\tan^{-1}n,\qquad r=90°-\theta_{\mathrm B} \]
\phtoku
電卓の $\tan^{-1}$ で $\tan^{-1}1.50\approx56.3°$。$r=90°-56.3°=33.7°$。
\phtashikame
屈折の法則で確かめる。
\[ \sin56.3°\approx0.832,\qquad 1.50\times\sin33.7°\approx1.50\times0.555\approx0.832 \]
で、合っている。極端な場合：$n\to1$（ガラスがない）なら $45°$ に近づく。
\kotae 入射角 約 $56°$、屈折角 約 $34°$
\end{phreidai}
\begin{jibun}
水（屈折率1.33）の表面のブルースター角を求めましょう。
\bkjans{$\tan^{-1}1.33\approx53.1°$。約 $53°$}
\end{jibun}

\begin{phreidai}[干渉計で鏡の動きをはかる]
\label{ph:波-13:ex12}波長 \qty{633}{nm} のレーザーを使ったマイケルソン干渉計で、鏡1をゆっくり動かすと、しまが1000本動きました。鏡はどれだけ動きましたか。
\phbamen
図\ref{ph:波-13:fig-michelson}の干渉計です。工作機械の台の位置をはかるのと同じしくみです。
\phmushi
空気の屈折率のゆらぎ（温度・気圧の変化）、鏡の傾きを無視します。
\phzu
しまの本数 $M=1000$、$\lambda=633\times10^{-9}$\,m、鏡の動き $\Delta$。
\phriyu
鏡が $\Delta$ 動くと、往復の経路は $2\Delta$ 変わる。経路の差が $\lambda$ 変わるごとにしまが1本動くので、⑦の $\Delta=M\lambda/2$。
\phtateru
\[ \Delta=\frac{M\lambda}{2} \]
\phtoku
\[ \Delta=\frac{1000\times633\times10^{-9}}{2}=3.165\times10^{-4}\ \mathrm{m}\approx\qty{0.32}{mm} \]
\phtashikame
単位：m。しま1本は $\lambda/2\approx\qty{0.32}{\micro m}$、その1000倍で \qty{0.32}{mm}。極端な場合：鏡が止まっていれば、しまは動かない（$M=0$）。
\kotae 約 \qty{0.32}{mm}
\end{phreidai}
\begin{jibun}
緑のレーザー（\qty{532}{nm}）を使った干渉計で、鏡を \qty{0.10}{mm} 動かすと、しまは何本動きますか。
\bkjans{$M=2\Delta/\lambda=2\times0.10\times10^{-3}\div(532\times10^{-9})$ $\approx376$。約380本}
\end{jibun}

\begin{phreidai}[表計算で夕焼けを作る]
\label{ph:波-13:ex13}空気の分子だけが光を散らすとします。太陽が真上のとき、空気の層全体を通る光は、波長 \qty{550}{nm} で $e^{-0.097}$ 倍（約91\,\%）に減ります。ほかの波長では、指数の 0.097 が $1/\lambda^4$ に比例して変わります。太陽が真上（空気の量1倍）、高さ $30°$（2倍）、地平線（38倍）のとき、\qty{450}{nm}・\qty{550}{nm}・\qty{700}{nm} の光が届く割合を、表計算で求めましょう。
\phbamen
図\ref{ph:波-13:fig-yuuyake}の、昼と夕方の太陽の光です。
\phmushi
ちり・水の粒・オゾンによる散乱や吸収、空気による光の曲がり、散らされた光がもう一度こちらへ来ることを無視します。地平線の38倍は、地球の丸さを入れて計算した値です。
\phzu
波長 $\lambda$、指数の数 $\tau=0.097\times(550/\lambda)^4$（\textbf{光学的な厚さ}）、空気の量の倍率 $X$。届く割合 $e^{-\tau X}$。
\phriyu
同じ厚さの空気を通るごとに、光は同じ割合ずつ散らされて減る。「同じ割合で減る」ものは指数関数になる（6-4③の $u=u_0e^{-kt}$ と同じ形）。空気の量が $X$ 倍なら、指数も $X$ 倍。太陽の高さ $h$ のとき、地球の丸さを無視すれば $X\approx1/\sin h$（$h=30°$ で2）。
\phtateru
\[ \text{届く割合}=e^{-\tau X},\qquad \tau=0.097\times\Bigl(\frac{550}{\lambda}\Bigr)^4 \]
\phtoku
\begin{center}
\bksheet{5}{19}{%
1 & $\lambda$（nm） & $\tau$ & 真上 $X=1$ & $30°$ $X=2$ & 地平線 $X=38$ \\
2 & 450 & 0.2165 & 0.805 & 0.649 & 0.00027 \\
3 & 550 & 0.0970 & 0.908 & 0.824 & 0.025 \\
4 & 700 & 0.0370 & 0.964 & 0.929 & 0.245 \\}
\end{center}
B2 は \texttt{=0.097*(550/A2)\^{}4}、C2 は \texttt{=EXP(-B2*1)}、D2 は \texttt{=EXP(-B2*2)}、E2 は \texttt{=EXP(-B2*38)}。下へコピーする。
\phtashikame
真上では、青も8割残るので、昼の太陽は白っぽく見える。地平線では、赤が約25\,\%残るのに、青は約0.03\,\%しか残らず、赤は青の約900倍。だから夕日は赤い。極端な場合：$X=0$（空気がない）なら、どの色も1（宇宙で見る太陽は白い）。
\kotae 真上：青 約81\,\%・緑 約91\,\%・赤 約96\,\%。地平線：青 約0.03\,\%・緑 約2.5\,\%・赤 約25\,\%
\end{phreidai}
\begin{jibun}
太陽の高さが約 $6°$ で、空気の量が10倍のとき、\qty{450}{nm}・\qty{550}{nm}・\qty{700}{nm} の光が届く割合を求めましょう。
\bkjans{$e^{-0.2165\times10}\approx0.11$、$e^{-0.097\times10}\approx0.38$、$e^{-0.037\times10}\approx0.69$。青 約11\,\%、緑 約38\,\%、赤 約69\,\%}
\end{jibun}

%% ---------------- まちがえやすいところ ----------------
\bksection{tsumazuki}{まちがえやすいところをチェック}
\begin{tsumazuki}{ヤングの式の $d$ に、すきまのはばを入れる}
\bkng{はば \qty{0.05}{mm} のすきまが2つ、中心の間が \qty{0.20}{mm} なので、$d=\qty{0.05}{mm}$。}
\bkok{$d$ は2つのすきまの中心の間かくで、$d=\qty{0.20}{mm}$。}
$\Delta x=L\lambda/d$ の $d$ は、2つの光の経路の差 $d\sin\theta$ を作る長さです。すきまのはば $a$ は、1つのすきまの中の回折（⑤）で、明るい線の明るさの変わり方に効きます。また、$L$・$\lambda$・$d$ は、必ず同じ単位（m）にそろえてから計算します。
\end{tsumazuki}
\begin{tsumazuki}{膜の中の波長を忘れる}
\bkng{厚さ \qty{300}{nm} の水の膜で強め合うのは、$2t=(m+\tfrac12)\lambda$ から \qty{400}{nm}。}
\bkok{膜の中では波長が $\lambda/n$ になるので、光路長の差は $2nt$。$2nt=(m+\tfrac12)\lambda$ から \qty{532}{nm}。}
物質の中の長さは、屈折率をかけた光路長で数えます。空気の層（くさび・ニュートンリング）だけは、$n=1$ なので $2t$ でかまいません。
\end{tsumazuki}
\begin{tsumazuki}{反射の入れかわりを、いつも1つと決める}
\bkng{反射防止膜も、シャボン膜と同じく「弱め合う：$2nt=m\lambda$」で厚さを決める。}
\bkok{反射防止膜は、表の面（空気→膜）でも裏の面（膜→ガラス）でも入れかわる。2つは打ち消し合い、弱め合うのは $2nt=(m+\tfrac12)\lambda$。}
入れかわりが起きるのは、屈折率の\textbf{大きい}物質との境で反射するときだけです。膜の両側の屈折率を見て、境ごとに1つずつ確かめましょう（表\ref{ph:波-13:tab-hansha}）。
\end{tsumazuki}
\begin{tsumazuki}{自然光にも、マリュスの法則を使う}
\bkng{自然光を偏光板1枚に通すと、$\cos^2\theta$ 倍になるので、向きによっては全部通る。}
\bkok{自然光は、どの向きに偏光板を回しても $\tfrac12$ になる。}
マリュスの法則の $\theta$ は、「来る光のゆれの向き」と「通す向き」の角です。自然光にはゆれの向きが決まっていないので、$\theta$ が決まりません。あらゆる向きの $\cos^2$ を平均して $\tfrac12$ になります。
\end{tsumazuki}

%% ---------------- 工学ではこう使う ----------------
\bksection{phgear}{工学ではこう使う}
\begin{phkougaku}[光の波で、分ける・はかる・守る]
\phboxhead{① 分光分析}
回折格子の分光器は、工場の品質検査、環境の測定（空気や水にふくまれる物質）、医療の検査、天文学で使われています。物質ごとに、吸ったり出したりする光の波長が決まっているので、光を波長ごとに分けて強さをはかれば、何がどれだけふくまれているかがわかります。スマホのカメラの前に回折格子のシートを置くだけで、かんたんな分光器が作れます。

\phboxhead{② 光で長さをはかる}
レーザー干渉計は、工作機械や半導体の製造装置で、台の位置を \unit{nm} の細かさではかります。ニュートンリングやオプチカルフラットの干渉じまは、レンズや鏡の面の形の検査に使われます。干渉計の考えは、光ファイバを使ったセンサー（ひずみや温度をはかる）にも使われています。

\phboxhead{③ 薄い膜で、光をえらぶ}
反射防止膜のほか、膜を何十層も重ねて、決まった波長だけを反射し、ほかは通す鏡（ダイクロイックミラー）が作られています。プロジェクターの色分け、レーザーの鏡、窓ガラスの「熱線反射」（赤外線だけを反射する膜。熱-13 の放射）などに使われています。

\phboxhead{④ 偏光で、見えない力を見る}
液晶の画面のほかに、偏光は材料の検査にも使われます。透明な樹脂で作った部品の模型に力をかけ、2枚の偏光板の間で見ると、応力（力-8）の大きい所に色のしまが現れます（光弾性）。また、偏光カメラは、ガラスや水面の反射を消したり、表面の傷を見つけたりするのに使われます。

\begin{phketa}[光の波の数]
\phketaline{可視光の波長}{約 \qtyrange{380}{780}{nm}}
\phketaline{赤いレーザー（ヘリウムネオン）}{\qty{633}{nm}}
\phketaline{CD・DVD・BD のトラックの間かく}{\qty{1.6}{\micro m}・\qty{0.74}{\micro m}・\qty{0.32}{\micro m}}
\phketaline{反射防止膜の厚さ}{約 \qty{100}{nm}}
\phketaline{目の分解能}{約1分（$3\times10^{-4}$\,rad）}
\phboxhead{答えが桁で合っているか}
干渉じまや光の点の位置は、「波長÷（間かくや直径）」の角で決まります。可視光は約 \qty{0.5}{\micro m} なので、間かくが \unit{\micro m} なら角は数十度、\unit{mm} なら $10^{-3}$\,rad（\qty{1}{m} 先で数 mm）が目安です。膜の厚さは、波長の4分の1から数倍（\qtyrange{0.1}{1}{\micro m}）のときに色が見えます。
\end{phketa}
\end{phkougaku}
\phstation{光学,分光分析,精密計測,光通信}{干渉、回折格子、分光器、光路長、薄膜、反射防止膜、分解能、開口数、偏光、ブルースター角、干渉計、散乱}

\begin{column}{色素のない色：構造色}
南アメリカのモルフォチョウの羽は、あざやかな青に光ります。ところが、この青は色素の色ではありません。羽の表面には、光の波長くらいの細かい層のつくりが並んでいて、そこで反射した光が干渉し、青い光だけが強め合うのです（④）。このように、物の細かい形（構造）が作る色を\textbf{構造色}といいます。\bkidx{こうぞうしょく}{構造色}クジャクの羽、タマムシ、シャボン玉、CD の色も構造色です。色素の色は、光に当たり続けると色あせますが、構造色は形が変わらないかぎり色あせません。そのため、構造色を使った、色あせない塗料や布の研究が進められています。
\end{column}

\begin{phsaki}{光は電磁波、そして粒}
この章では、光が横波であること、ゆれる電子が光を出すこと、反射で山と谷が入れかわることを、結果として使いました。これらはすべて、光が電場と磁場の波（電磁波）であることから導けます。物理シリーズの『電磁気学』（電-15）で、電場と磁場の式から光の波を導き、加速度運動をする電気が光を出すわけ、反射率の式（図\ref{ph:波-13:fig-fresnel}）を学びます。

一方、光は、エネルギーのとびとびの粒（光子）としてもふるまいます（原-2）。干渉じまを、光子を1つずつ数えられるほど弱い光で作ると、光子は1つずつスクリーンのばらばらの所に届き、たくさん集まると干渉じまになります。光は「波か粒か」ではなく、その両方の性質をもつのです。光学の専門の本では、回折をフーリエ変換で扱うフーリエ光学、レンズの像のゆがみ（収差）、ホログラム、何十層もの薄膜の設計を学びます。
\end{phsaki}

%% ---------------- 章末問題 ----------------
\bkendproblems
\bkproblemgroup{基本}{3}
\bkproblem 空気の分子による散乱で、紫の光（\qty{400}{nm}）は、だいだい色の光（\qty{600}{nm}）の何倍強く散らされますか（\nobreak\bkreidairef{1}\nobreak）。
\bkpans{$(600\div400)^4=1.5^4\approx5.1$。約5.1倍}
\bkproblem ナトリウムランプの光（\qty{589}{nm}）で、間かく \qty{0.50}{mm} の2つのすきまを使い、\qty{2.0}{m} 離れたスクリーンでヤングの実験をします。明るい線の間かくを求めましょう（\nobreak\bkreidairef{2}\nobreak）。
\bkpans{$\Delta x=2.0\times589\times10^{-9}$ $\div(0.50\times10^{-3})$ $\approx2.4\times10^{-3}$\,m。約 \qty{2.4}{mm}}
\bkproblem 間かく \qty{0.30}{mm} の2つのすきまと、\qty{1.5}{m} 離れたスクリーンでヤングの実験をすると、明るい線の間かくが \qty{2.7}{mm} でした。光の波長を求めましょう（\nobreak\bkreidairef{2}\nobreak）。
\bkpans{$\lambda=\Delta x\,d/L=2.7\times10^{-3}\times0.30\times10^{-3}\div1.5$ $=5.4\times10^{-7}$\,m。\qty{540}{nm}}
\bkproblem 1 mm あたり500本の回折格子に、\qty{650}{nm} の光を垂直に当てます。何次の明るい線まで現れますか。1次と2次の向きも求めましょう（\nobreak\bkreidairef{3}\nobreak）。
\bkpans{$d=2.0\times10^{-6}$\,m。$m\le d/\lambda\approx3.1$ なので3次まで。1次 $\sin\theta=0.325$ で約 $19°$、2次 $\sin\theta=0.65$ で約 $41°$（3次は約 $77°$）}
\bkproblem 水（屈折率1.33）の上に浮いた、厚さ \qty{250}{nm} の油の膜（屈折率1.45）を、ほぼ垂直に見ます。強く反射する可視光の波長を求めましょう（\nobreak\bkreidairef{5}\nobreak）。
\bkpans{表（空気→油）で入れかわり、裏（油→水、$1.45\to1.33$）はそのまま。強め合う $2nt=(m+\tfrac12)\lambda$、$2nt=725$\,nm。$m=1$ で $725\div1.5\approx\qty{483}{nm}$（約 \qty{480}{nm}、青）。$m=0$ は \qty{1450}{nm}、$m=2$ は \qty{290}{nm} で見えない}
\bkproblem フッ化マグネシウム（屈折率1.38）の反射防止膜で、\qty{500}{nm} の光の反射をいちばん弱める、いちばん薄い膜の厚さを求めましょう（\nobreak\bkreidairef{6}\nobreak）。
\bkpans{$t=500\div(4\times1.38)\approx\qty{91}{nm}$}
\bkproblem 長さ \qty{10}{cm} の2枚のガラス板のはしに、太さ \qty{0.030}{mm} の糸をはさみ、\qty{600}{nm} の光を上から当てます。暗い線の間かくを求めましょう（\nobreak\bkreidairef{7}\nobreak）。
\bkpans{$\Delta x=0.10\times600\times10^{-9}$ $\div(2\times0.030\times10^{-3})$ $=1.0\times10^{-3}$\,m。\qty{1.0}{mm}}
\bkproblem 球面の半径 \qty{2.0}{m} の平凸レンズで、\qty{589}{nm} の光のニュートンリングを見ます。5番めの暗い輪の半径を求めましょう（\nobreak\bkreidairef{8}\nobreak）。
\bkpans{$\rho_5=\sqrt{5\times589\times10^{-9}\times2.0}$ $\approx2.4\times10^{-3}$\,m。約 \qty{2.4}{mm}}
\bkproblemgroup{チャレンジ}{5}
\bkproblem 通す向きが直角の2枚の偏光板の間に、1枚めから角 $\theta$ の向きの偏光板をはさみます。自然光（強さ $I_0$）を通したとき、出てくる光の強さを $\theta$ の式で表し、いちばん明るくなる $\theta$ と、そのときの強さを求めましょう（\nobreak\bkreidairef{10}\nobreak・4-9④）。
\bkpans{$I=\tfrac12I_0\cos^2\theta\cos^2(90°-\theta)$ $=\tfrac12I_0\cos^2\theta\sin^2\theta$ $=\tfrac18I_0\sin^22\theta$。$\theta=45°$ でいちばん明るく、$\tfrac18I_0$（12.5\,\%）}
\bkproblem 夜、向こうから来る車の2つのライト（間かく \qty{1.5}{m}）を、2つの光として見分けられるのは、何 km までですか。ひとみの直径を \qty{3.0}{mm}、光の波長を \qty{550}{nm} とします。実際には、それより近くでないと見分けられないことが多いのはなぜですか（\nobreak\bkreidairef{9}\nobreak）。
\bkpans{$\theta_{\min}\approx2.24\times10^{-4}$\,rad。$1.5\div(2.24\times10^{-4})$ $\approx6.7\times10^3$\,m。約 \qty{6.7}{km}。夜はひとみが広がるが、空気のゆらぎ、目のレンズのゆがみ、網膜の細胞の間かくなどで、回折の限界まで見えないから}
\bkproblem マイケルソン干渉計の一方の腕に、長さ \qty{5.0}{cm} の、中を真空にした透明な容器を置きます。容器にゆっくり空気（1気圧・\qty{20}{\celsius}）を入れていくと、\qty{633}{nm} の光のしまが43本動きました。空気の屈折率を求めましょう。光は容器の中を往復します（\nobreak\bkreidairef{12}\nobreak・④の光路長）。
\bkpans{光路長の差の変化は往復で $2(n-1)L$。これが $M\lambda$ なので $n-1=M\lambda/(2L)$ $=43\times633\times10^{-9}\div(2\times0.050)$ $\approx2.7\times10^{-4}$。$n\approx1.00027$}
\bkproblem \bkreidairef{13}のモデルで、\qty{450}{nm} の青い光の届く割合が、\qty{700}{nm} の赤い光の届く割合のちょうど半分になるのは、空気の量が何倍のときですか。そのときの太陽の高さは、およそ何度ですか（$X\approx1/\sin h$ とする。5-14 の自然対数を使う）。
\bkpans{$e^{-0.2165X}\div e^{-0.0370X}=e^{-0.1795X}=\tfrac12$。$0.1795X=\ln2\approx0.693$ より $X\approx3.9$。$\sin h\approx1/3.86\approx0.259$ で、$h\approx15°$}

\begin{furikaeri}
\begin{checklist}
\item ヤングの実験・回折格子・薄い膜で、光が強め合う所と弱め合う所を、経路の差から計算できる。\bkfurref{ph:波-13:sec2}
\item 円い穴の回折で見分けられる角の限界と、偏光板を通る光の強さ・ブルースター角を計算できる。\bkfurref{ph:波-13:sec5}
\item ゆすられる電子のモデル（強制振動）から、空が青く夕焼けが赤いわけを説明できる。\bkfurref{ph:波-13:sec8}
\end{checklist}
\end{furikaeri}
\bkfukidashi{sora}{光が波だとわかると、色も、見える限界も、空の青さも、計算できたね。次の波-14 では、波を使って情報を送るしくみを見ていくよ。}
